
\begin{enumerate}
\def\labelenumi{\arabic{enumi})}
\tightlist
\item
  \begin{enumerate}
  \def\labelenumii{\alph{enumii})}
  \tightlist
  \item
    Let \(\omega, p, q \in \mathbf{R}\) with \(p \leqslant q\) . We set \(f(x) = e^{i\omega x}\) for \(p \leqslant x \leqslant q\) , \(f(x) = 0\) for \(x > q\) or \(x < p\) . Then for \(y \in \mathbf{R}\) , one has
  \end{enumerate}
\end{enumerate}

\[
\mathcal {F} (f) (y) = i \frac {e ^ {i p (\omega - 2 \pi y)} - e ^ {i q (\omega - 2 \pi y)}}{\omega - 2 \pi y}.
\]

\begin{enumerate}
\def\labelenumi{\alph{enumi})}
\setcounter{enumi}{1}
\tightlist
\item
  Let \(\omega \in \mathbf{R}\) , \(\beta > 0\) . We set \(f(x) = e^{(-\beta + i\omega)x}\) for \(x \geqslant 0\) , \(f(x) = 0\) for \(x < 0\) . Then for \(y \in \mathbf{R}\) , one has
\end{enumerate}

\[
\mathcal {F} (f) (y) = \frac {i}{\omega - 2 \pi y + i \beta}.
\]

\begin{enumerate}
\def\labelenumi{\alph{enumi})}
\setcounter{enumi}{2}
\tightlist
\item
  Let \(a > 0\) . We set \(f(x) = \frac{1}{a} e^{-a|x|}\) for \(x \in \mathbf{R}\) . Then for \(y \in \mathbf{R}\) , one has
\end{enumerate}

\[
\mathcal {F} (f) (y) = \frac {2}{4 \pi^ {2} y ^ {2} + a ^ {2}}.
\]

\(d)\) We set \(f(x) = e^{-\pi x^2}\) for \(x \in \mathbf{R}\) . Then for \(y \in \mathbf{R}\) , one has \(\mathcal{F}(f)(y) = e^{-\pi y^2}\) . (Find a first-order linear differential equation satisfied by the Fourier transform of \(f\) , and use the formula \(\int_0^{+\infty} e^{-t^2} dt = \frac{1}{2} \sqrt{\pi}\) of FVR, VII, §1, no. 3.)

\begin{enumerate}
\def\labelenumi{\alph{enumi})}
\setcounter{enumi}{4}
\tightlist
\item
  Let \(n \in N\) and let Q be a positive definite quadratic form on \(R^{n}\) . Let \(\sigma\) a linear map on \(R^{n}\) such that \(\mathrm{Q}(x) = \| \sigma(x) \|^{2}\) . Let \(Q^{*}\) the positive definite quadratic form such that \(\mathrm{Q}^{*}(y) = \|^{t} \sigma^{-1}(y) \|^{2}\) for all \(y \in R^{n}\) . Let \(z \in C\) a complex number such that \(\mathcal{I}(z) > 0\) . The function f from \(R^{n}\) to C defined by \(f(x) = \exp(i\pi z\mathrm{Q}(x))\) belongs to \(\mathcal{S}(\mathbf{R}^{n})\) and we have
\end{enumerate}

\[
\mathscr {F} (f) (y) = | \det (\sigma) | ^ {- 1} \left(\frac {i}{z}\right) ^ {n / 2} \exp (i \pi z ^ {- 1} Q ^ {*} (y))
\]

for all \(y \in R^{n}\) .

\(f)\) Let \(f \in \mathcal{C}(\mathbf{R})\) the function defined by \(f(x) = (\sin(\pi x)/(\pi x))^2\) for \(x \neq 0\) and \(f(0) = 1\) . We have \(\mathcal{F}(f)(y) = 0\) if \(|y| > 1\) and \(\mathcal{F}(f)(y) = 1 - |y|\) otherwise. Let \(g(x) = xf(x)\) . We have \(\mathcal{F}(g)(y) = 0\) if \(|y| > 1\) and \(\mathcal{F}(g)(y) = s(y)/(2i\pi)\) otherwise.

\begin{enumerate}
\def\labelenumi{\arabic{enumi})}
\setcounter{enumi}{1}
\tightlist
\item
  We identify the interval {[}0,1{[} with \(\mathbf{T} = \mathbf{R} / \mathbf{Z}\) by the canonical projection. Let \(0 \leqslant a < b < 1\) . Let \(f\) the continuous function from \(\mathbf{R} / \mathbf{Z}\) into \(\mathbf{R}\) such that \(f(x) = 0\) if \(x \notin [a,b]\) , \(f((a + b) / 2) = 1\) , and such that \(f\) is affine on the intervals \([a, (a + b) / 2]\) and \([(a + b) / 2, b]\) .
\end{enumerate}

Compute the Fourier transform of \(f\) and prove that there exists a constant \(C \geqslant 0\) such that

\[
| \widehat {f} (n) | \leqslant \inf \left(b - a, \frac {\mathrm{C}}{n ^ {2}}\right)
\]

for all \(n \in \mathbf{Z} - \{0\}\) .

\begin{enumerate}
\def\labelenumi{\arabic{enumi})}
\setcounter{enumi}{2}
\tightlist
\item
  \begin{enumerate}
  \def\labelenumii{\alph{enumii})}
  \tightlist
  \item
    The functions \(t \mapsto t^n e^{-\pi t^2}\) for \(n\) integer \(\geqslant 0\) form a total set in \(L^2(\mathbf{R})\) . (Let \(f \in L^2(\mathbf{R})\) an element orthogonal to the \(t^n e^{-\pi t^2}\) . For \(z \in \mathbf{C}\) , set \(F(z) = \int_{-\infty}^{+\infty} f(t)e^{-\pi t^2}e^{2i\pi tz}dt\) , and show that \(F\) is an entire function whose derivatives at 0 are all zero. Deduce that the Fourier transform of the function \(t \mapsto f(t)e^{-\pi t^2}\) on \(\mathbf{R}\) is zero.)
  \end{enumerate}
\end{enumerate}

\begin{enumerate}
\def\labelenumi{\alph{enumi})}
\setcounter{enumi}{1}
\item
  For \(k \in N\) , one has \(\partial_{t}^{k}(e^{-2\pi t^{2}}) = \mathrm{P}_{k}(t)e^{-2\pi t^{2}}\) , where \(P_{k}\) is a polynomial of degree k, with leading coefficient \((-4\pi)^{k}\) .
\item
  Let \(k \in N\) , \(P \in C[t]\) nonzero and \(\alpha\) the leading coefficient of P. Then
\end{enumerate}

\[
\int_ {\mathbf {R}} \mathrm{P} _ {k} (t) e ^ {- 2 \pi t ^ {2}} \mathrm{P} (t)   d t = \left\{ \begin{array}{l l} 0 & \text {if} \deg (\mathrm{P}) <   k \\ (- 1) ^ {k} k! 2 ^ {- 1 / 2} \alpha & \text {if} \deg (\mathrm{P}) = k. \end{array} \right.
\]

\begin{enumerate}
\def\labelenumi{\alph{enumi})}
\setcounter{enumi}{3}
\tightlist
\item
  For \(k \in \mathbf{N}\) and \(t \in \mathbf{R}\) , we set:
\end{enumerate}

\[
\mathcal {H} _ {k} (t) = (- 1) ^ {k} \frac {2 ^ {1 / 4 - k}}{\sqrt {\pi^ {k} k !}} e ^ {\pi t ^ {2}} \partial_ {t} ^ {k} (e ^ {- 2 \pi t ^ {2}})
\]

(“Hermite functions”). The family \((\mathcal{H}_{k})\) is an orthonormal basis of the space \(\mathrm{L}^{2}(\mathbf{R})\) . (Use a), b), c)).

\begin{enumerate}
\def\labelenumi{\alph{enumi})}
\setcounter{enumi}{4}
\tightlist
\item
  We have \(\mathcal{F}(\mathcal{H}_k) = (-i)^k\mathcal{H}_k\)
\end{enumerate}

\begin{enumerate}
\def\labelenumi{\arabic{enumi})}
\setcounter{enumi}{3}
\tightlist
\item
  \begin{enumerate}
  \def\labelenumii{\alph{enumii})}
  \tightlist
  \item
    Let \(\alpha > 0\) , \(h > 0\) , and \(f: \mathbf{R} \to \mathbf{R}\) the continuous function vanishing outside the interval \([- \alpha, \alpha[\) , linear in \([- \alpha, 0]\) and \([0, \alpha]\) , equal to \(h\) at 0. We have:
  \end{enumerate}
\end{enumerate}

\[
\mathcal {F} (f) (t) = \frac {h}{\alpha \pi^ {2}} \frac {\sin^ {2} (\pi \alpha t)}{t ^ {2}}
\]

for \(t \in R^{*}\) .

\begin{enumerate}
\def\labelenumi{\alph{enumi})}
\setcounter{enumi}{1}
\tightlist
\item
  Let \(\mathbf{N} \in \mathbf{N}\) . Set \(g(x) = \sum_{n=-\mathbf{N}}^{\mathbf{N}} f(x + n)\) . Then:
\end{enumerate}

\[
\mathcal {F} (g) (t) = \frac {h}{\alpha \pi^ {2}} \frac {\sin^ {2} (\pi \alpha t)}{t ^ {2}} \frac {\sin ((2 N + 1) \pi t)}{\sin (\pi t)}
\]

for \(t\in \mathbf{R}^*\)

\begin{enumerate}
\def\labelenumi{\alph{enumi})}
\setcounter{enumi}{2}
\item
  We choose sequences \((\alpha_{i})_{i\geqslant1}\) , \((h_{i})_{\geqslant1}\) , \((\mathrm{N}_{i})_{i\geqslant1}\) , whence, for each integer i, a function \(g_{i}\) constructed by the procedure of b). If \(\sum_{i} h_{i} < +\infty\) , the series \(\sum_{i} g_{i}\) converges uniformly to a continuous function G on R.
\item
  If
\end{enumerate}

\[
\sum_ {i} \alpha_ {i} h _ {i} \mathrm{N} _ {i} <   + \infty , \quad \sum_ {i} \frac {h _ {i}}{\alpha_ {i}} \log \mathrm{N} _ {i} <   + \infty , \quad \sum_ {i} h _ {i} \mathrm{N} _ {i} = + \infty ,
\]

then G and \(\mathcal{F}(\mathrm{G})\) are integrable on \(\mathbf{R}\) but

\[
\sum_ {n = - \infty} ^ {+ \infty} \mathrm{G} (n) = + \infty .
\]

Show that there exist examples of such sequences.

\begin{enumerate}
\def\labelenumi{\arabic{enumi})}
\setcounter{enumi}{4}
\item
  Let \(\mu \in \mathcal{M}^1 (\mathrm{G})\) . We assume that \(\mathcal{F}(\mu)\in \mathrm{L}^1 (\widehat{\mathrm{G}})\) . Then \(\mu\) is a measure with base the Haar measure on G (INT, V, § 5, no. 2, def. 2), and admits a continuous density with respect to it. (Use the filter base \(\mathfrak{B}\) of prop. 8 of II, p.~211. Show that \(\lim_{\varphi ,\mathfrak{B}}\varphi *\mu = \mu\) for the vague topology, and for \(\varphi\) in an element of \(\mathfrak{B}\) , one has \(\mathcal{F}(\varphi *\mu) = \mathcal{F}(\varphi)\mathcal{F}(\mu)\) ; but \(\mathcal{F}(\mu)\in \mathrm{L}^1 (\widehat{\mathrm{G}})\) , hence \(\varphi *\mu = \overline{\mathcal{F}} (\mathcal{F}(\varphi *\mu))\) tends to \(\overline{\mathcal{F}} (\mathcal{F}(\mu))\) into \(\mathcal{C}_0(\mathrm{G})\) . Thus \(d\mu (x) = f(x)dx\) where \(f = \overline{\mathcal{F}} (\mathcal{F}(\mu))\) .
\item
  Let H be a subgroup of G, and \(\mu \in \mathcal{M}^{1}(\widehat{\mathrm{G}})\) . For \(\mathcal{F}(\mu)\) to be invariant under translations by elements of H, it is necessary and sufficient that the support of \(\mu\) be contained in \(H^{\perp}\) .
\item
  Let \(f \in \mathrm{L}^1(\mathrm{G})\) . Let \(\mathrm{F}\) a holomorphic function defined in an open set \(\mathrm{U}\) of \(\mathbf{C}\) . We assume that \(\mathrm{U}\) contains the set of values of \(\mathcal{F}(f)\) , and if \(\mathrm{G}\) is not discrete, we assume that \(0 \in \mathrm{U}\) and \(\mathrm{F}(0) = 0\) .
\end{enumerate}

Then there exists \(g \in \mathrm{L}^{1}(\mathrm{G})\) such that \(\mathrm{F} \circ \mathcal{F}(f) = \mathcal{F}(g)\) (Use the holomorphic functional calculus.)

\begin{enumerate}
\def\labelenumi{\arabic{enumi})}
\setcounter{enumi}{7}
\tightlist
\item
  Let \(p \in [1,2]\) and \(q \in [2, +\infty]\) such that \(1/p + 1/q = 1\) .
\end{enumerate}

\begin{enumerate}
\def\labelenumi{\alph{enumi})}
\item
  Let \(f \in \mathcal{K}(G)\) . Show that \(\|\mathcal{F}(f)\|_{q} \leqslant \|f\|_{p}\) (This inequality is valid for p = 1 and p = 2. In the general case, use the M. Riesz inequality (INT, IV, §6, exerc. 18).)
\item
  Deduce that the restriction of \(\mathcal{F}\) to \(\mathcal{K}(\mathrm{G})\) extends to a continuous linear map from \(\mathrm{L}^p (\mathrm{G})\) into \(\mathrm{L}^q (\widehat{\mathrm{G}})\) which coincides with the Fourier transform on \(\mathrm{L}^p (\mathrm{G})\cap \mathrm{L}^1 (\mathrm{G})\) and on \(\mathrm{L}^p (\mathrm{G})\cap \mathrm{L}^2 (\mathrm{G})\) .
\end{enumerate}

9) Let G be a locally compact group (not necessarily commutative) equipped with a left Haar measure.

\begin{enumerate}
\def\labelenumi{\alph{enumi})}
\item
  If \(f \in \mathrm{L}^{1}(\mathrm{G})\) , there exists \(f_{1}\) and \(f_{2}\) into \(\mathrm{L}^{1}(\mathrm{G})\) such that \(f_{1} * f_{2} = f\) (Use exercise 30 of I, p.~185.)
\item
  Let \(f_{1}, f_{2} \in \mathrm{L}^{1}(\mathrm{G})\) with \(f_{1} \geqslant 0, f_{2} \geqslant 0\) . Show that \(f_{1} * f_{2}\) is almost everywhere equal to a lower semicontinuous function. (Consider \(f_{1}\) as the simple limit of an increasing sequence of functions \(\geqslant 0\) of \(\mathrm{L}^{\infty}(\mathrm{G})\) .)
\item
  Take G = R. For every open interval I = {]}a, b{[}, set
\end{enumerate}

\[
\mathrm{A} (\mathrm{I}) = ] a, a + (b - a) / 6 [, \quad \mathrm{A} ^ {\prime} (\mathrm{I}) = ] a + 5 (b - a) / 6, b [.
\]

Define intervals \(I_{n}\) for \(n \geqslant 1\) and \(I_{n}^{\prime}\) for \(n \geqslant 0\) by \(I_{0}^{\prime} = ]0, 1[\) and

\[
\mathrm{I} _ {n} = \mathrm{A} (\mathrm{I} _ {n - 1} ^ {\prime}), \qquad \mathrm{I} _ {n} ^ {\prime} = \mathrm{A} ^ {\prime} (\mathrm{I} _ {n - 1} ^ {\prime})
\]

for \(n \geqslant 1\) . Let F be the union of the intervals \(\mathrm{I}_n\) for \(n \geqslant 1\) . There exists no pair \((\mathrm{F}_1, \mathrm{F}_2)\) of subsets of \(\mathbf{R}\) such that \(\mathrm{F}_1 + \mathrm{F}_2 = \mathrm{F}\) .

\begin{enumerate}
\def\labelenumi{\alph{enumi})}
\setcounter{enumi}{3}
\tightlist
\item
  If f is a continuous function \(\geqslant 0\) on R such that \(F = \{x \in \mathbf{R} | f(x) > 0\}\) , then there exists no pair of functions \((f_{1}, f_{2})\) into \(L^{1}(\mathbf{R})\) such that \(f_{1} \geqslant 0\) , \(f_{2} \geqslant 0\) and \(f = f_{1} * f_{2}\) . (Suppose such an equality. Let \(E_{i}\) the set of \(x \in R\) such that \(f_{i}(x) > 0\) . Let \(F_{i}\) be the set of points of density of \(E_{i}\) (INT, V, §6, exerc. 15). Prove that \(F = F_{1} + F_{2}\) .)
\end{enumerate}

\begin{enumerate}
\def\labelenumi{\arabic{enumi})}
\setcounter{enumi}{9}
\item
  Let \(k \in Z\) . The map \(x \mapsto x^{k}\) the map from G to G is injective if and only if the map \(\chi \mapsto \chi^{k}\) of \(\widehat{G}\) into \(\widehat{G}\) is surjective.
\item
  Let G be a finite abelian group and d an integer. \(\geqslant 1\) Prove that for all \(x \in G\) and every \(\chi \in \widehat{G}\) , one has
\end{enumerate}

\[
\sum_{\substack{y\in \mathrm{G}\\ y^{d} = x}}\chi (y) = \sum_{\substack{\eta \in \widehat{\mathrm{G}}\\ \eta^{d} = \chi}}\eta (x).
\]

\begin{enumerate}
\def\labelenumi{\arabic{enumi})}
\setcounter{enumi}{11}
\item
  Let \((\mathrm{G}_{i})_{i\in\mathrm{I}}\) a family of locally compact abelian groups. For every \(i\in I\) , let \(H_{i}\) an open compact subgroup of \(G_{i}\) . Let G be the local direct product of the \(G_{i}\) with respect to the \(H_{i}\) For every \(\chi\in\widehat{G}\) , let \(\chi_{i}\) the restriction of \(\chi\) to \(G_{i}\) . Show that \(\chi\mapsto(\chi_{i})_{i\in\mathrm{I}}\) is an isomorphism of the topological group \(\widehat{G}\) onto the local direct product of the \(\widehat{G}_{i}\) with respect to the \(H_{i}^{\perp}\) .
\item
  Let \(G_{d}\) the group G equipped with the discrete topology and \(\chi\) a unitary character of \(G_{d}\) . For every integer \(n \geqslant 1\) , all elements \(x_{1}, \ldots, x_{n} \in G\) and every \(\varepsilon > 0\) , there exists \(\widehat{x} \in \widehat{G}\) such that \(|\langle\widehat{x}, x_{i}\rangle - \chi(x_{i})| \leqslant \varepsilon\) for \(1 \leqslant i \leqslant n\) . (The canonical injection of \(G_{d}\) into G has as its dual a morphism from \(\widehat{G}\) into \(\widehat{G}_{d}\) whose image is dense by Cor. 5 of II, p.~229.)
\item
  \begin{enumerate}
  \def\labelenumii{\alph{enumii})}
  \tightlist
  \item
    The space \(\mathrm{L}^{1}(\mathrm{G})\) is a closed ideal in \(\mathcal{M}^{1}(\mathrm{G})\) , hence \(\widehat{G}\) identifies with an open subset of \(\mathsf{X}(\mathcal{M}^{1}(\mathrm{G}))\) (cf. exercise 17 of I, p. 169).
  \end{enumerate}
\end{enumerate}

\begin{enumerate}
\def\labelenumi{\alph{enumi})}
\setcounter{enumi}{1}
\tightlist
\item
  Let \(\mathcal{M}_{d}(\mathrm{G})\) (resp. \(\mathcal{M}_{a}(\mathrm{G})\) ) the set of \(\mu\in\mathcal{M}^{1}(\mathrm{G})\) which are diffuse (resp. atomic). Then \(\mathcal{M}_{d}(\mathrm{G})\) is a closed ideal of \(\mathcal{M}^{1}(\mathrm{G})\) and \(\mathcal{M}_{a}(\mathrm{G})\) is a closed subalgebra of \(\mathcal{M}^{1}(\mathrm{G})\) .
\end{enumerate}

15) For every integer \(k \geqslant 1\) for all elements \(x_1, \ldots, x_k\) pairwise distinct in E and for all integers \(n_1, \ldots, n_k\) , the relation \(x_1^{n_1} \cdots x_k^{n_k} = e\) implies \(x_1^{n_1} = \cdots = x_k^{n_k} = e\) . We say that E is a Kronecker set if every continuous complex-valued function on E with absolute value 1 is a uniform limit on E of unitary characters of G. If \(q\) is an integer \(\geqslant 1\) and G = (Z/qZ) \(^{\mathbf{N}}\) we say that E is of type K \(_q\) if every continuous map from E into the set of roots \(q\) characters of the unit coincide on E with a unitary character of G.

\begin{enumerate}
\def\labelenumi{\alph{enumi})}
\item
  Let P be a compact Kronecker set. If \(\mu \in \mathcal{M}^1 (\mathrm{G})\) is concentrated on P, we have \(\| \mathcal{F}(\mu)\|_{\infty} = \| \mu \|\) Hence deduce that every complex function continuous on P is the restriction of the Fourier transform of an element. \(f\in \mathrm{L}^{1}(\widehat{\mathrm{G}})\) .
\item
  Let \(E \subset G\) a finite independent set. Let f be a complex function of absolute value 1 on E. Suppose that \(f(x)^{q} = 1\) when \(x \in E\) is of order q. The function f is the uniform limit on E of unitary characters of G. (Use exerc. 13.)
\item
  Every Kronecker set is independent and contains only elements of infinite order.
\item
  Every set of type \(K_{q}\) is independent.
\item
  Let \(k \geqslant 1\) an integer and let \(V_{1}, \ldots, V_{k}\) of nonempty disjoint open subsets of G. If every neighborhood of e in G contains an element of infinite order, there exists \(x_{i} \in V_{i}\) for \(1 \leqslant i \leqslant k\) such that \(\{x_{1}, \ldots, x_{k}\}\) let be a Kronecker set. If \(\mathrm{G} = (\mathbf{Z}/q\mathbf{Z})^{\mathbf{N}}\) , there exists \(x_{i} \in V_{i}\) such that \(\{x_{1}, \ldots, x_{k}\}\) let be of type \(K_{q}\) .
\item
  Let \(E \subset G\) a compact independent set, \(F = E \cup E^{-1}\) . Let \(\mu\) a bounded diffuse measure on G, concentrated on F. Then the measures \(\mu^{n}\) , for \(n \geqslant 0\) are pairwise disjoint, where one sets \(\mu^{0} = \varepsilon_{e}\) and \(\mu^{n} = \mu * \mu^{n-1}\) for \(n \geqslant 1\) . (Show that, if m \textless{} n, the set of \((x_{1}, \ldots, x_{n}) \in G^{n}\) such that \(x_{1} \ldots x_{n} \in F^{m}\) is negligible for \(\mu \otimes \cdots \otimes \mu\) ). Deduce that, if \(\mu \geqslant 0\) , one has \(\| \sum_{k=0}^{n} \alpha_{k} \mu^{k} \| = \sum_{k=0}^{n} |\alpha_{k}| \| \mu \|^{k}\) whatever the complex numbers \(\alpha_{0}, \ldots, \alpha_{n}\) .
\end{enumerate}

\(g)\) Let \(r \geqslant 1\) an integer. Let E \(\subset\) G a compact set independent, \(\mu_1, \ldots, \mu_r\) of measures \(\geqslant 0\) of mass 1, concentrated on disjoint parts E \(_1, \ldots, \text{E}_r\) of E. Let \(z_1, \ldots, z_r \in \mathbf{C}\) with \(|z_i| \leqslant 1\) for all \(i\) . There exists a character \(\chi\) of \(\mathcal{M}^1(\text{G})\) such that \(\chi(\mu_i) = z_i\) for \(1 \leqslant i \leqslant r\) (First show, by reasoning as in \(f\) ), that if \((n_1, \ldots, n_r) \neq (m_1, \ldots, m_r)\) the measures \(\mu_1^{m_1} * \cdots * \mu_r^{m_r}\) and \(\mu_1^{n_1} * \cdots * \mu_r^{n_r}\) are disjoint. Next, assuming \(|z_i| = 1\) for all \(i\) , show that the spectral radius of \(\varepsilon_e + \overline{z}_1 \mu_1 + \cdots + \overline{z}_r \mu_r\) is \(r+1\) , whence a character \(\chi\) of \(\mathcal{M}^1(\text{G})\) such that \(\chi(\varepsilon_e + \overline{z}_1 \mu_1 + \cdots + \overline{z}_r \mu_r) = r+1\) ; this character answers the question. In the general case, write \(z_i = \frac{1}{2}(z_i' + z_i'')\) with \(|z_i'| = |z_i''| = 1\) , and \(\mu_i = \frac{1}{2}(\mu_i' + \mu_i'')\) where \(\mu_1', \mu_1'', \ldots, \mu_r', \mu_r''\) satisfy the same hypotheses as \(\mu_1, \ldots, \mu_r\) .)

\begin{enumerate}
\def\labelenumi{\alph{enumi})}
\setcounter{enumi}{7}
\tightlist
\item
  Let \(\mathcal{M}_{d}(\mathrm{E})\) the Banach space of diffuse measures on G concentrated on E. Every continuous linear form of norm \(\leqslant 1\) on \(\mathcal{M}_{d}(\mathrm{E})\) extends to a character of \(\mathcal{M}^{1}(\mathrm{G})\) (Use g).
\end{enumerate}

\begin{enumerate}
\def\labelenumi{\arabic{enumi})}
\setcounter{enumi}{15}
\tightlist
\item
  For every complex number \(s\) such that \(\mathcal{R}(s) > 1\) , we denote by
\end{enumerate}

\[
\zeta (s) = \sum_ {n \geqslant 1} \frac {1}{n ^ {s}}
\]

(“Riemann zeta function”).

\begin{enumerate}
\def\labelenumi{\alph{enumi})}
\item
  The function \(\zeta\) is holomorphic on the open set of \(\mathbf{C}\) consisting of complex numbers \(s\) such that \(\mathcal{R}(s) > 1\) .
\item
  For \(y\in \mathbf{R}_+^*\) , let
\end{enumerate}

\[
\theta (y) = \sum_ {n \in {\bf Z}} e ^ {- \pi n ^ {2} y};
\]

This series converges absolutely and uniformly on compact subsets of \(R_{+}^{*}\) . For every s such that \(\mathcal{R}(s)>1\) , one has

\[
\pi^ {- s / 2} \Gamma (s / 2) \zeta (s) = \frac {1}{2} \int_ {0} ^ {+ \infty} \theta (y) y ^ {s / 2} \frac {d y}{y}
\]

(see FVR, VII, p.~10 for the gamma function in the complex plane).

\begin{enumerate}
\def\labelenumi{\alph{enumi})}
\setcounter{enumi}{2}
\item
  For every \(y \in R_{+}^{*}\) , one has \(\theta(y^{-1}) = \sqrt{y}\theta(y)\) . (Use the Poisson formula and exercise 1.)
\item
  For every \(s\) such that \(\mathcal{R}(s) > 1\) , one has
\end{enumerate}

\[
\pi^ {- s / 2} \Gamma (s / 2) \zeta (s) = \frac {1}{s (s - 1)} + \frac {1}{2} \int_ {1} ^ {+ \infty} (\theta (y) - 1) (y ^ {s / 2} + y ^ {(1 - s) / 2}) \frac {d y}{y}.
\]

\begin{enumerate}
\def\labelenumi{\alph{enumi})}
\setcounter{enumi}{4}
\item
  There exists a unique holomorphic function on \(\mathbf{C} = \{1\}\) whose restriction to the open set of complex numbers with real part \(>1\) coincides with \(\zeta\) . This function, still denoted \(\zeta\) , satisfies \((s - 1)\zeta(s) \to 1\) when \(s \to 1\) (*in other words, it has a simple pole with residue 1 at \(s = 1_*\) ).
\item
  The function \(\Lambda\) defined by
\end{enumerate}

\[
\Lambda (s) = \pi^ {- s / 2} \Gamma (s / 2) \zeta (s)
\]

is holomorphic on \(\mathbf{C} = \{0,1\}\) ; it has simple poles with residue 1 at \(s = 0\) and \(s = 1\) . We have \(\Lambda(1 - s) = \Lambda(s)\) for all \(s \in \mathbf{C} = \{0,1\}\) .

\(g)\) Show that for all \(s \in \mathbf{C}\) such that \(\mathcal{R}(s) > 1\) , one has

\[
\zeta (s) = \prod_ {p} (1 - p ^ {- s}) ^ {- 1},
\]

where the product is taken over the set of prime numbers and is absolutely convergent (“Euler product”).

\begin{enumerate}
\def\labelenumi{\alph{enumi})}
\setcounter{enumi}{7}
\tightlist
\item
  When \(\sigma > 1\) converges to 1, we have
\end{enumerate}

\[
\sum_ {p} p ^ {- \sigma} = \log \left(\frac {1}{\sigma - 1}\right) + O (1),
\]

where the sum is over the prime numbers.

\begin{enumerate}
\def\labelenumi{\roman{enumi})}
\tightlist
\item
  Everything \(s \in C - \{1\}\) such that \(\zeta(s) = 0\) is either of the form \(s = -2k\) , where \(k \geqslant 1\) is an integer, or satisfies \(0 \leqslant \mathcal{R}(s) \leqslant 1\) .
\end{enumerate}

\begin{enumerate}
\def\labelenumi{\alph{enumi})}
\setcounter{enumi}{9}
\tightlist
\item
  We have \(\zeta(s) \neq 0\) if \(\mathcal{R}(s) = 1\) and \(s \neq 1\) . (Show using the Euler product that
\end{enumerate}

\[
\zeta (\sigma) ^ {3} \zeta (\sigma + i t) ^ {4} \zeta (\sigma + 2 i t) \geqslant 1
\]

for all \(\sigma > 1\) and \(t \in \mathbf{R}\) , then consider \(t\) such that \(\zeta(1 + it) = 0\) .

\begin{enumerate}
\def\labelenumi{\arabic{enumi})}
\setcounter{enumi}{16}
\tightlist
\item
  Let \(q \geqslant 1\) an integer, and let \(a \geqslant 1\) be an integer coprime to \(q\) . Let \(\widetilde{\chi}\) a character of the group \((\mathbf{Z}/q\mathbf{Z})^*\) . One denotes by \(\chi\) the map from \(\mathbf{Z}\) into \(\mathbf{C}\) such that \(\chi(n) = 0\) if \(n\) is not coprime to \(q\) and \(\chi(n) = \widetilde{\chi}(n \pmod{q})\) in the opposite case (“Dirichlet characters”).
\end{enumerate}

\begin{enumerate}
\def\labelenumi{\alph{enumi})}
\tightlist
\item
  Suppose \(\widetilde{\chi}\) nontrivial. The function defined by
\end{enumerate}

\[
\mathrm{L} (s, \chi) = \sum_ {n \geqslant 1} \chi (n) n ^ {- s}
\]

is holomorphic on the open set U of C consisting of the complex numbers such that \(\mathcal{R}(s) > 1\) and there exists a unique holomorphic function on \(\mathcal{R}(s) > 0\) which coincides with it on U. We have

\[
\mathrm{L} (s, \chi) = \prod_ {p} (1 - \chi (p) p ^ {- s}) ^ {- 1}
\]

for \(s \in U\) , where the product is over the prime numbers and converges absolutely and uniformly on compact subsets of U (“Dirichlet L-function”).

\begin{enumerate}
\def\labelenumi{\alph{enumi})}
\setcounter{enumi}{1}
\item
  Suppose that \(\widetilde{\chi}^{2}\) is not trivial. We have \(L(1,\chi)\neq0\) . (Adapt the method of question j) of the previous exercise.)
\item
  Suppose that \(\widetilde{\chi}\) has order 2. For every integer \(n \geqslant 1\) , we denote by
\end{enumerate}

\[
r _ {\chi} (n) = \sum_ {d | n} \chi (d)
\]

where the sum is over the divisors \(d \geqslant 1\) of n. There exists a real number c \textgreater{} 0 such that, for every \(x \geqslant 2\) , one has

\[
\sum_ {n \leqslant x} \frac {r _ {\chi} (n)}{\sqrt {n}} \geqslant c \log (x)
\]

(note that the map \(r_{\chi}\) is \(\geqslant 0\) is multiplicative, that is, \(r_{\chi}(nm) = r_{\chi}(n)r_{\chi}(m)\) if n and m are coprime, and note that \(r_{\chi}(n) \geqslant 1\) if n is the square of an integer.)

\(d)\) One still assumes that \(\widetilde{\chi}\) has order 2. One has

\[
\sum_ {n \leqslant x} \frac {r _ {\chi} (n)}{\sqrt {n}} = \frac {1}{2} \mathrm{L} (1, \chi) x ^ {1 / 2} + \mathrm{O} (1)
\]

when \(x \to +\infty\) . Deduce that \(\mathrm{L}(1, \chi) \neq 0\) .

\begin{enumerate}
\def\labelenumi{\alph{enumi})}
\setcounter{enumi}{4}
\tightlist
\item
  When \(\sigma > 1\) converges to 1, one has
\end{enumerate}

\[
\sum_ {p \equiv a \text {mod.} q} p ^ {- \sigma} = \frac {1}{\varphi (q)} \sum_ {p} p ^ {- \sigma} + \mathrm{O} (1),
\]

where \(\varphi(q)\) is the Euler indicator of \(q\) (A, V, p. 76).

\begin{enumerate}
\def\labelenumi{\alph{enumi})}
\setcounter{enumi}{5}
\tightlist
\item
  There exist infinitely many primes congruent to a modulo q (“Dirichlet's arithmetic progression theorem”).
\end{enumerate}

\begin{enumerate}
\def\labelenumi{\arabic{enumi})}
\setcounter{enumi}{17}
\tightlist
\item
  \begin{enumerate}
  \def\labelenumii{\alph{enumii})}
  \tightlist
  \item
    Let \(n \in N\) and let \(f \in \mathcal{S}(\mathbf{R}^{n})\) . The function of \(R \times R^{n}\) to C defined by
  \end{enumerate}
\end{enumerate}

\[
u (t, x) = \int_ {\mathbf {R} ^ {n}} \widehat {f} (y) \exp (2 \pi i x \cdot y - 4 \pi^ {2} t \| y \| ^ {2}) d y
\]

for \((t,x)\in\mathbf{R}\times\mathbf{R}^{n}\) satisfies

\[
\left\{ \begin{array}{l} u (0, x) = f (0, x) \\ \partial_ {t} u = i \Delta_ {2} u, \end{array} \right.\tag{1}
\]

(“Schrödinger equation”) where

\[
\Delta_ {2} u = \sum_ {i = 1} ^ {n} \partial_ {x _ {i}} ^ {2} u.
\]

\begin{enumerate}
\def\labelenumi{\alph{enumi})}
\setcounter{enumi}{1}
\tightlist
\item
  The function \(u\) is the unique solution of (1) such that \(t \mapsto (x \mapsto u(t,x))\) belongs to \(\mathcal{C}^{\infty}(\mathbf{R},\mathcal{S}(\mathbf{R}^{n}))\) .
\end{enumerate}

\begin{enumerate}
\def\labelenumi{\arabic{enumi})}
\setcounter{enumi}{18}
\item
  Let \(n \in N\) and let \(f \in \mathrm{L}^{1}(\mathbf{R}^{n})\) (resp. \(\mathrm{L}^{2}(\mathbf{R}^{n})\) ) such that \(f \circ \sigma = f\) for every element \(\sigma\) of the special orthogonal group \(\mathrm{SO}(n, \mathbf{R})\) . Then the Fourier transform \(\widehat{f}\) of f satisfies \(\widehat{f}(\sigma(y)) = \widehat{f}(y)\) for all \(\sigma \in \mathrm{SO}(n, \mathbf{R})\) and every \(y \in R^{n}\) (resp. almost every \(y \in R^{n}\) ).
\item
  Let \(n \geqslant 1\) an integer. One identifies \(R_{+}^{*} \times S_{n}\) , and \(R^{n+1} - \{0\}\) by the map \((t, z) \mapsto tz\) . Let \(\omega\) the unique measure on the sphere \(S_{n} \subset R^{n+1}\) invariant under \(\mathbf{O}(n + 1, \mathbf{R})\) such that the measure \(t^{n}dt \otimes \mu\) is identified with the Lebesgue measure on \(R^{n+1} - \{0\}\) (cf. INT, VII, p. 116, exercise 8). Let \(\mu\) be the measure on \(R^{n+1}\) image of the measure \(\omega\) by the injection of \(S_{n}\) into \(R^{n+1}\) .
\end{enumerate}

\begin{enumerate}
\def\labelenumi{\alph{enumi})}
\item
  We have \(\mu (\mathbf{S}_n) = 2\pi^{(n + 1) / 2}\Gamma (\frac{n + 1}{2})^{-1}\) (cf. INT, V, §8, no. 7).
\item
  For every real number \(k \geqslant 0\) , one defines the function \(J_{k} \colon R_{+} \to C\) by
\end{enumerate}

\[
\mathrm{J} _ {k} (t) = \frac {(t / 2) ^ {k}}{\Gamma (k + \frac {1}{2}) \Gamma (\frac {1}{2})} \int_ {- 1} ^ {1} e ^ {i u t} (1 - u ^ {2}) ^ {k - 1 / 2} d u
\]

for \(t \in R_{+}\) (“Bessel function of the first kind”). For every \(x \in R^{n+1}\) , one has

\[
\mathcal {F} (\mu) (x) = \frac {2 \pi}{\| x \| ^ {(n - 1) / 2}} \mathrm{J} _ {(n - 1) / 2} (2 \pi \| x \|)
\]

where \(\|x\|\) is the Euclidean norm.

\begin{enumerate}
\def\labelenumi{\alph{enumi})}
\setcounter{enumi}{2}
\tightlist
\item
  If \(k\in \mathbf{N}\) , one has
\end{enumerate}

\[
\mathrm{J} _ {k} (t) = \frac {1}{2 \pi} \int_ {0} ^ {2 \pi} \cos (t \sin (u) - k u) d u
\]

for all \(t \in R_{+}\) .

\(d)\) If \(n\) is even, show that \(\mathcal{F}(\mu)(x)\) is a polynomial in \(\| x\|^{\pm 1/2}\) , \(\cos(2\pi \| x\|)\) and \(\sin(2\pi \| x\|)\) ; for \(n=2\) , one has

\[
\mathcal {F} (\mu) (x) = \frac {\sin (2 \pi \| x \|)}{2 \pi \| x \|}.
\]

\begin{enumerate}
\def\labelenumi{\alph{enumi})}
\setcounter{enumi}{4}
\tightlist
\item
  Let g be a continuous function with compact support in \(R_{+}^{*}\) . Set \(f(x) = g(\|x\|^{2}) \|x\|^{1-n/2}\) for all \(x \in R^{n}\) . Then we have \(\widehat{f}(y) = h(\|y\|^{2}) \|y\|^{1-n/2}\) for \(y \in R^{n}\) , where
\end{enumerate}

\[
h (s) = \pi \int_ {0} ^ {+ \infty} g (t) \mathrm{J} _ {n / 2 - 1} (2 \pi \sqrt {t s}) d t
\]

for all \(s \in \mathbf{R}_{+}\) (“Hankel transform”).

21) Let H be a closed subgroup of G and f an element of L¹(G) such that \(\mathcal{F}(f)\) vanishes on H⊥. Let \(\varepsilon > 0\) . There exists a measure \(\mu \in \mathcal{M}^1(G)\) concentrated on H such that \(\| \mu \| \leqslant 2\) , \(\| f * \mu \| \leqslant \varepsilon\) and such that \(\mathcal{F}(\mu) = 1\) in a neighbourhood of H⊥.

\begin{enumerate}
\def\labelenumi{\arabic{enumi})}
\setcounter{enumi}{21}
\tightlist
\item
  Let \(p\) a prime number. Let \(K\) a finite extension of the field \(\mathbf{Q}_p\) of \(p\) -adic numbers; it is a locally compact non-discrete field. Let \(t: K \to \mathbf{Q}_p\) the trace map. Denote by \(\lambda: \mathbf{Q}_p \to \mathbf{Q}\) the map defined in prop. 20 of II, p.~237.
\end{enumerate}

\begin{enumerate}
\def\labelenumi{\alph{enumi})}
\tightlist
\item
  The additive group of K is in duality with itself by means of the map \((x,y)\mapsto\exp(2i\pi\lambda(t(xy)))\) . One identifies K and its dual by this map.
\end{enumerate}

For every integer \(n \geqslant 1\) , we denote by \(\mathcal{S}(\mathrm{K}^{n})\) the complex vector space of functions from \(K^{n}\) to C that are locally constant and have compact support.

\begin{enumerate}
\def\labelenumi{\alph{enumi})}
\setcounter{enumi}{1}
\item
  The restriction of the Fourier transform to \(\mathcal{S}(K^{n})\) is an automorphism whose inverse is the restriction of the Fourier cotransform.
\item
  For every real number \(p \in [1, +\infty[\) , the space \(\mathcal{S}(\mathrm{K}^n)\) is dense in \(\mathrm{L}^p (\mathrm{K}^n)\) .
\end{enumerate}

\begin{enumerate}
\def\labelenumi{\arabic{enumi})}
\setcounter{enumi}{22}
\tightlist
\item
  Let \(p \geqslant 3\) a prime number and \(k \geqslant 1\) an integer. We denote \(G = F_{p}^{k}\) . One equips G with the Haar measure of total mass 1, and one identifies \(\widehat{G}\) and G via the map that to \(n = (n_{1}, \ldots, n_{k}) \in G\) associates the character
\end{enumerate}

\[
e _ {n} \colon x \mapsto \exp \left(\frac {2 i \pi}{p} \sum_ {j = 1} ^ {k} n _ {j} x _ {j}\right).
\]

For functions \(f_{1}, f_{2}, f_{3}\) from G to \(\mathbf{C}\) , we denote by

\[
\Lambda (f _ {1}, f _ {2}, f _ {3}) = \frac {1}{p ^ {2 k}} \sum_ {x \in \mathrm{G}} \sum_ {h \in \mathrm{G}} f _ {1} (x) f _ {2} (x + h) f _ {3} (x + 2 h).
\]

We denote \(\| f\|\) the norm of \(f\) into \(\mathrm{L}^2 (\mathrm{G})\) .

\begin{enumerate}
\def\labelenumi{\alph{enumi})}
\tightlist
\item
  We have
\end{enumerate}

\[
\Lambda (1, f _ {2}, f _ {3}) = \left(\frac {1}{p ^ {k}} \sum_ {x \in \mathrm{G}} f _ {2} (x)\right) \left(\frac {1}{p ^ {k}} \sum_ {x \in \mathrm{G}} f _ {3} (x)\right)
\]

and

\[
\Lambda (1, 1, f _ {3}) = \frac {1}{p ^ {k}} \sum_ {x \in \mathrm{G}} f _ {3} (x).
\]

\begin{enumerate}
\def\labelenumi{\alph{enumi})}
\setcounter{enumi}{1}
\tightlist
\item
  We have
\end{enumerate}

\[
\Lambda (f _ {1}, f _ {2}, f _ {3}) = \sum_ {t \in \mathrm{G}} \widehat {f} _ {1} (t) \widehat {f} _ {2} (- 2 t) \widehat {f} _ {3} (t),
\]

and

\[
| \Lambda (f _ {1}, f _ {2}, f _ {3}) | \leqslant \| f _ {1} \| \| f _ {2} \| \| \widehat {f _ {3}} \| _ {\infty}.
\]

\begin{enumerate}
\def\labelenumi{\alph{enumi})}
\setcounter{enumi}{2}
\tightlist
\item
  Let \(f \colon \mathbf{G} \to \mathbf{R}\) a function such that \(\sum f(x) = 0\) . There exists an affine hyperplane \(\mathrm{H} \subset \mathrm{G}\) such that
\end{enumerate}

\[
\frac {1}{p ^ {k - 1}} \sum_ {x \in \mathrm{H}} f (x) \geqslant \frac {1}{2} \| \widehat {f} \| _ {\infty}
\]

(show that there exist elements \(t \in \mathbf{G}\) and \(\theta \in \mathbf{R}\) such that the real part of \(p^{-k} \sum_{x} f(x)(e_t(x) \exp(i\theta) + 1)\) is equal to \(\| \widehat{f} \|_{\infty}\) ).

\begin{enumerate}
\def\labelenumi{\alph{enumi})}
\setcounter{enumi}{3}
\tightlist
\item
  Let A ⊂ G contain no arithmetic progression of length 3, that is, there do not exist x ∈ A and h ∈ G = \{0\} such that x + h and x + 2h also belong to A. Show that there exists an affine hyperplane H ⊂ G such that
\end{enumerate}

\[
\frac {\operatorname{Card} (\mathrm{A} \cap \mathrm{H})}{\operatorname{Card} (\mathrm{H})} \geqslant \frac {\operatorname{Card} (\mathrm{A})}{\operatorname{Card} (\mathrm{G})} + \frac {1}{2} \left(\left(\frac {\operatorname{Card} (\mathrm{A})}{\operatorname{Card} (\mathrm{G})}\right) ^ {2} - \frac {1}{p ^ {k}}\right)
\]

(apply the previous questions to \(f = \varphi_{A} - |A|/p^{k}\) , where \(\varphi_{A}\) is the characteristic function of A).

\begin{enumerate}
\def\labelenumi{\alph{enumi})}
\setcounter{enumi}{4}
\tightlist
\item
  Let A ⊂ G contain no arithmetic progression of length 3. Show that \(\text{Card}(A) \leqslant 3p^{k}/k\) ("Roth's theorem for \(F_{p}^{k}\) »).
\end{enumerate}

24) a) Let δ \textgreater{} 0. If p is a sufficiently large prime number and if A is a subset of Z/pZ such that Card(A) ≥ δp, then A contains an arithmetic progression of length 3. (Adapt the method of the preceding exercise.)

\begin{enumerate}
\def\labelenumi{\alph{enumi})}
\setcounter{enumi}{1}
\tightlist
\item
  Let \(\delta > 0\) . If N is a sufficiently large integer and if A \(\subset \{1, \ldots, N\}\) satisfies \(\operatorname{Card}(A) \geqslant \delta N\) then A contains an arithmetic progression of length 3. (Reduce to the case of the previous question by choosing \(p\) appropriately as a function of N.)
\end{enumerate}

\begin{enumerate}
\def\labelenumi{\arabic{enumi})}
\setcounter{enumi}{24}
\tightlist
\item
  We denote \(\mu_{G}\) the given Haar measure on G. Let \(\mu\) a complex measure on G.
\end{enumerate}

\begin{enumerate}
\def\labelenumi{\alph{enumi})}
\item
  The image of the Fourier transform of \(\mu\) is finite if and only if \(\mu\) is a linear combination of idempotent measures, that is, of measures \(\nu\) such that \(\nu * \nu = \nu\) .
\item
  The support of \(\mathcal{F}_{\mathrm{G}}(\mu)\) is finite if and only if \(\mu\) is a finite linear combination of measures \(\widehat{x}\mu_{G}\) , where \(\widehat{x} \in \widehat{G}\) .
\item
  We denote \(A_{\mu}\) the set of measures of the form \(\widehat{x}\mu\) for \(\widehat{x}\in\widehat{G}\) . The set of measures \(\nu\in A_{\mu}\) such that \(\nu(\mathrm{G})\neq0\) is finite if and only if there exists a closed subgroup H of G, a finite set K, characters \(\widehat{x}_{k}\) for \(k\in K\) and complex numbers \(t_{k}\) for \(k\in K\) such that
\end{enumerate}

\[
\nu = \Bigl (\sum_ {k \in \mathrm{K}} t _ {k} \widehat {x} _ {k} \Bigr) i (\mu_ {\mathrm{H}})
\]

where i denotes the canonical inclusion of H into G.

\begin{enumerate}
\def\labelenumi{\arabic{enumi})}
\setcounter{enumi}{25}
\tightlist
\item
  For every compact abelian group H, we denote by \(\mu_{H}\) the normalized Haar measure on H.
\end{enumerate}

Assume G is compact. We denote \(E_{G}\) the set of complex measures \(\mu\) on G such that the Fourier transform of \(\mu\) takes integer values. This is a commutative subgroup of \(\mathcal{M}^{1}(G)\) which contains the Haar measure \(\mu_{G}\) . One denotes by \(F_{G}\) the subgroup of \(E_{G}\) generated by the measures \(\widehat{x} i_{\mathrm{H}}(\mu_{\mathrm{H}})\) , where H ranges over the closed subgroups of G \(i_{H}\) denoting the canonical inclusion of H into G), and where \(\widehat{x}\) ranges over \(\widehat{G}\) .

The goal of this exercise is to prove that \(E_{G} = F_{G}\) (Cohen's idempotent theorem).

\begin{enumerate}
\def\labelenumi{\alph{enumi})}
\item
  Let \(\varphi\colon\mathrm{G}_{1}\to\mathrm{G}_{2}\) a morphism of compact commutative groups. The map \(\mu\mapsto\varphi(\mu)\) is a group homomorphism from \(\mathrm{E}_{\mathrm{G}_{1}}\) into \(\mathrm{E}_{\mathrm{G}_{2}}\) .
\item
  For every \(\mu \in \mathrm{E}_{\mathrm{G}}\) and every \(\widehat{x} \in \widehat{\mathrm{G}}\) , one has \(\widehat{x}\mu_{\mathrm{G}} \in \mathrm{E}_{\mathrm{G}}\) .
\end{enumerate}

Let \(\mu\) a complex measure on G. We denote \(\theta\) the function on G such that \(\mu = \theta |\mu|\) . One denotes by \(\mathrm{A}_{\mu} \subset \mathcal{M}^{1}(\mathrm{G})\) the set of measures of the form \(\widehat{x}\mu\) . Let \(\nu\) a point adherent to \(\mathrm{A}_{\mu}\) into \(\mathcal{M}^{1}(\mathrm{G})\) .

\begin{enumerate}
\def\labelenumi{\alph{enumi})}
\setcounter{enumi}{2}
\tightlist
\item
  If \(\nu\) does not belong to \(\mathrm{A}_{\mu}\) , then \(\nu\) is singular with respect to \(\mu_{\mathrm{G}}\) .\\
\item
  Let \(\varepsilon > 0\) . Let \(f \in \mathcal{C}(\mathrm{G})\) of norm \(\leqslant 1\) such that \(\nu(f) > (1 - \varepsilon)\|\mu\|\) and \(\widehat{x} \in \widehat{\mathrm{G}}\) such that \(\mathcal{R}(\mu(\widehat{x} f)) > (1 - \varepsilon)\|\mu\|\) . Show that
\end{enumerate}

\[
\int_ {\mathrm{G}} | 1 - f \widehat {x} \theta | d | \mu | \leqslant 3 \varepsilon^ {1 / 2} \| \mu \|.
\]

\begin{enumerate}
\def\labelenumi{\alph{enumi})}
\setcounter{enumi}{4}
\item
  Suppose that \(\mu \in \mathrm{E}_{\mathrm{G}}\) . Show that \(\| \nu \| \leqslant \| \mu \| - 1 / (36 \| \mu \|)\) . (Apply the preceding question to two characters \(\widehat{x}_1\) and \(\widehat{x}_2\) such that \(\widehat{x}_1 \mu \neq \widehat{x}_2 \mu\) , and note that \(\| \widehat{x}_1 \mu - \widehat{x}_2 \mu \| \geqslant 1\) .)
\item
  Conclude that \(E_{G} = F_{G}\) and more precisely that every measure \(\mu\) into \(E_{G}\) can be written as
\end{enumerate}

\[
\mu = \sum_ {k \in \mathrm{K}} n _ {k} \widehat {x} _ {k} i (\mu_ {\mathrm{H} _ {k}})
\]

where the set K is finite, \(n_{k} \in Z\) , \(\widehat{x}_{k} \in \widehat{G}\) , and the measures \(i(\mu_{\mathrm{H}_{k}})\) are mutually disjoint. (Let \(\mu \in E_{G}\) ; show that there exists a measure \(\nu \neq 0\) in the closure of \(\{\nu \in A_{\mu} \mid \nu(G) \neq 0\}\) whose norm is minimal; then show, using e), that \(\{\widehat{x} \in \widehat{G} \mid \nu(\widehat{x}) \neq 0\}\) is finite, and apply the preceding exercise.)

\begin{enumerate}
\def\labelenumi{\arabic{enumi})}
\setcounter{enumi}{26}
\tightlist
\item
  Let \(f\) a function in \(\mathcal{S}(\mathbf{R})\) of norm \(\mathrm{L}^2\) equal to 1. We have (uncertainty principle) the inequality
\end{enumerate}

\[
\left(\int_ {\mathbf {R}} x ^ {2} | f (x) | ^ {2} d x\right) \left(\int_ {\mathbf {R}} t ^ {2} | \widehat {f} (t) | ^ {2} d t\right) \geqslant \frac {1}{(4 \pi) ^ {2}}
\]

(observe that

\[
- \frac {1}{2} = \mathcal {R} \Bigl (\int_ {\mathbf {R}} x f (x) \overline {{f ^ {\prime} (x)}} d x \Bigr),
\]

and apply the Cauchy-Schwarz inequality). Determine the cases of equality.

\begin{enumerate}
\def\labelenumi{\arabic{enumi})}
\setcounter{enumi}{27}
\tightlist
\item
  \begin{enumerate}
  \def\labelenumii{\alph{enumii})}
  \tightlist
  \item
    Let \(\delta > 0\) a real number. There does not exist a function \(f: \mathbf{R} \to \mathbf{C}\) integrable, nonzero, and compactly supported, such that \(\widehat{f}(y) = \mathrm{O}(\exp(-\delta|y|))\) when \(|y| \to +\infty\) (In particular, there does not exist such a function such that \(\widehat{f}\) also compactly supported.)
  \end{enumerate}
\end{enumerate}

\begin{enumerate}
\def\labelenumi{\alph{enumi})}
\setcounter{enumi}{1}
\tightlist
\item
  Let \((\varrho_{n})_{n\geqslant1}\) a summable sequence of strictly positive real numbers. The product
\end{enumerate}

\[
g (x) = \prod_ {n \geqslant 1} \frac {\sin (\varrho_ {n} x)}{\varrho_ {n} x},
\]

defined as equal to 1 for x = 0, converges absolutely and uniformly on every compact subset of R. It defines a continuous, nonzero, even function such that \(g \in \mathrm{L}^{1}(\mathbf{R})\) .

\begin{enumerate}
\def\labelenumi{\alph{enumi})}
\setcounter{enumi}{2}
\tightlist
\item
  Let \(\varrho\) the sum of the series \(\sum \varrho_{n}\) The support of the Fourier transform of \(g\) is contained in \([- \varrho, \varrho]\) .
\end{enumerate}

\(d)\) Let \(r > 0\) a real number. Let \(\delta: \mathbf{R}_{+} \to \mathbf{R}_{+}^{*}\) a decreasing function such that \(\delta(y) \to 0\) when \(y \to +\infty\) , such that \(y \mapsto \delta(y)/y\) is integrable on \([1, +\infty[\) with respect to Lebesgue measure, and such that \(y\delta(y) \to +\infty\) when \(y \to +\infty\) . There exists a continuous function \(f \in \mathrm{L}^{1}(\mathbf{R})\) with compact support contained in \([-r, r]\) such that

\[
\widehat {f} (y) = \mathrm{O} (\exp (- \delta (| y |) | y |))
\]

when \(|y| \to +\infty\) (First show that there exists a summable decreasing sequence \((\varrho_n)\) of strictly positive real numbers such that \(\varrho \leqslant r\) and \(\varrho_n \geqslant e\delta(n)n^{-1}\) for all \(n\) large enough.)

\begin{enumerate}
\def\labelenumi{\arabic{enumi})}
\setcounter{enumi}{28}
\tightlist
\item
  Let \(f\in \mathrm{L}^2 (\mathbf{R})\) such that
\end{enumerate}

\[
\int_ {\mathbf {R}} | f (x) | ^ {2} e ^ {x ^ {2}} d x \leqslant 1, \quad \int_ {\mathbf {R}} | \widehat {f} (y) | ^ {2} e ^ {y ^ {2}} d y \leqslant 1.
\]

\begin{enumerate}
\def\labelenumi{\alph{enumi})}
\tightlist
\item
  For every integer \(n \in \mathbf{N}\) , the function \(g_n\) defined by
\end{enumerate}

\[
g _ {n} (z) = \left(\int_ {\mathbf {R}} f (x) x ^ {n} \exp (- z x ^ {2} / 2) d x\right) ^ {2}
\]

is holomorphic in the open set \(\mathrm{U} = \{z \in \mathbf{C} \mid \mathcal{R}(z) > -1\}\) of \(\mathbf{C}\) .

\begin{enumerate}
\def\labelenumi{\alph{enumi})}
\setcounter{enumi}{1}
\tightlist
\item
  For every integer \(n \in \mathbf{N}\) , the function \(h_n\) defined by
\end{enumerate}

\[
h _ {n} (z) = \frac {(- 1) ^ {n}}{z ^ {2 n + 1}} \left(\int_ {\mathbf {R}} \widehat {f} (y) y ^ {n} \exp (- y ^ {2} / (2 z)) d y\right) ^ {2}
\]

is holomorphic in the open set \(\mathrm{V} = \{z \in \mathbf{C} \mid \mathcal{R}(z^{-1}) > -1\} = \{z \in \mathbf{C} \mid |z + 1/2| > 1/2\}\) of \(\mathbf{C}\) .

\begin{enumerate}
\def\labelenumi{\alph{enumi})}
\setcounter{enumi}{2}
\item
  There exists a holomorphic function \(k_{0}\) into \(U \cup V = C - \{-1\}\) such that \(k_{0}|U = g_{0}\) and \(k_{0}|V = h_{0}\) .
\item
  The function \(k_{0}\) is zero. (Show that there exists \(c \in C\) such that \(|k_{0}(z)| \leqslant c/|z+1|\) for all \(z \in C - \{-1\}\) and apply Liouville's theorem to \(z \mapsto (z+1)k_{0}(z).\) )
\item
    Prove by induction on \(n \in \mathbf{N}\) that there exists a holomorphic function \(k_n\) into \(\mathrm{U} \cup \mathrm{V} = \mathbf{C} - \{-1\}\) such that \(k_n | \mathrm{U} = g_n\) and \(k_n | \mathrm{V} = h_n\) , then that \(k_n = 0\) f) We have \(f = 0\) .
\item
  Let \(a, b > 0\) such that \(ab \geqslant 1\) . If \(f \in \mathrm{L}^2(\mathbf{R})\) satisfies
\end{enumerate}

\[
\int_ {\mathbf {R}} | f (x) | ^ {2} e ^ {a x ^ {2}} d x <   + \infty , \qquad \int_ {\mathbf {R}} | \widehat {f} (y) | ^ {2} e ^ {b y ^ {2}} d y <   + \infty ,
\]

then \(f = 0\) (“Hardy’s uncertainty principle”).

\begin{enumerate}
\def\labelenumi{\arabic{enumi})}
\setcounter{enumi}{29}
\tightlist
\item
  \begin{enumerate}
  \def\labelenumii{\alph{enumii})}
  \tightlist
  \item
    For every \(x_0 \in \mathbf{T}\) , there exists \(f \in \mathcal{C}(\mathbf{T})\) such that the Fourier series of \(f\) to \(x_0\) does not converge symmetrically to \(f(x_0)\) . (Let \(s_n(f)\) the symmetric partial sums of the Fourier series of \(f\) ; consider the linear maps \(f \mapsto s_n(f)(x_0)\) and apply the Banach–Steinhaus theorem.)
  \end{enumerate}
\end{enumerate}

\begin{enumerate}
\def\labelenumi{\alph{enumi})}
\setcounter{enumi}{1}
\item
  There exists \(f \in \mathrm{L}^{1}(\mathbf{T})\) such that the Fourier series of f does not converge symmetrically to f in \(\mathrm{L}^{1}(\mathbf{T})\) .
\item
  The Fourier transform is not surjective from \(\mathrm{L}^{1}(\mathbf{T})\) in the Banach space of sequences \((c_{n})_{n\in\mathbf{Z}}\) tending to 0 at infinity. (Apply Cor. 1 of EVT, I, p.~19.)
\end{enumerate}

\(d)\) Suppose that \(f \in \mathcal{C}^1(\mathbf{T})\) . The sequence \(s_n(f)\) converges to \(f\) into \(\mathcal{C}(\mathbf{T})\) .

\begin{enumerate}
\def\labelenumi{\arabic{enumi})}
\setcounter{enumi}{30}
\tightlist
\item
  Let \(r \geqslant 1\) an integer. We denote \(\mathcal{S}(\mathbf{T}^r)\) the complex vector space of infinitely differentiable functions on \(\mathbf{T}^r\) , and \(\mathcal{S}(\mathbf{Z}^r)\) the complex vector space of functions \(f\colon \mathbf{Z}^r \to \mathbf{C}\) such that, for every real \(A > 0\) , one has \(f(n) = O(\|n\|^{-A})\) when \(\|n\| \to +\infty\) .
\end{enumerate}

\begin{enumerate}
\def\labelenumi{\alph{enumi})}
\tightlist
\item
  Equipped with the seminorms
\end{enumerate}

\[
p _ {\mathrm{A}} (f) = \sup _ {n \in \mathbf {Z} ^ {r}} \| n \| ^ {\mathrm{A}} | f (n) |,
\]

the space \(\mathcal{S}(\mathbf{Z}^{r})\) is a Fréchet space.

\begin{enumerate}
\def\labelenumi{\alph{enumi})}
\setcounter{enumi}{1}
\item
  For every real number \(p \in [1, +\infty[\) , the space \(\mathcal{S}(\mathbf{Z}^{r})\) (resp. the space \(\mathcal{S}(\mathbf{T}^{r})\) ) is dense in \(\mathrm{L}^{p}(\mathbf{Z}^{r})\) (respectively in \(\mathrm{L}^{p}(\mathbf{T}^{r})\) ).
\item
  The restriction of the Fourier transform of \(T^{r}\) to \(\mathcal{S}(\mathbf{T}^{r})\) is an isomorphism from \(\mathcal{S}(\mathbf{T}^{r})\) into \(\mathcal{S}(\mathbf{Z}^{r})\) , whose inverse is the restriction of the Fourier cotransform.
\end{enumerate}

\(d)\) The topological vector space \(\mathcal{S}(\mathbf{Z}^{r})\) is isomorphic to \(\mathcal{S}(\mathbf{R}^{n})\) , in particular \(c\) is a Montel space. (Use Exercise 3.)

\begin{enumerate}
\def\labelenumi{\arabic{enumi})}
\setcounter{enumi}{31}
\tightlist
\item
  Let G be a finite abelian group. We endow G with the Haar measure of total mass 1.
\end{enumerate}

\begin{enumerate}
\def\labelenumi{\alph{enumi})}
\tightlist
\item
  Show that if f is nonzero, then
\end{enumerate}

\[
\operatorname{Card} (\operatorname{Supp} (f)) \operatorname{Card} (\operatorname{Supp} (\widehat {f})) \geqslant \operatorname{Card} (\mathrm{G})
\]

Uncertainty principle; note that

\[
\| \widehat {f} \| _ {\infty} \leqslant \frac {1}{\mathrm{Card(G)}} \sum_ {x \in \mathrm{G}} | f (x) |
\]

and apply the Cauchy-Schwarz inequality and the Plancherel formula.)

\begin{enumerate}
\def\labelenumi{\alph{enumi})}
\setcounter{enumi}{1}
\tightlist
\item
  Determine the cases of equality.
\end{enumerate}

\begin{enumerate}
\def\labelenumi{\arabic{enumi})}
\setcounter{enumi}{32}
\tightlist
\item
  For every integer \(r \geqslant 1\) , we denote by \(\mathrm{V}_{r} = \det((x_{i}^{j})_{\substack{1 \leqslant i \leqslant r \\ 0 \leqslant j < r}}) \in \mathbf{Z}[x_{1}, \ldots, x_{r}]\) the Vandermonde determinant (A, III, p.~99, example 1).
\end{enumerate}

\begin{enumerate}
\def\labelenumi{\alph{enumi})}
\tightlist
\item
  Let \(f_{1}, \ldots, f_{r}\) infinitely differentiable functions on R. For all real numbers \(x_{1} < \cdots < x_{r}\) , there exist real numbers \((t_{0}, \ldots, t_{r-1})\) into \([x_{1}, x_{r}]\) such that
\end{enumerate}

\[
\det (f_{i}(x_{j})) = \frac{1}{1!2!\cdots(r - 1)!}\mathrm{V}_{r}(x_{1},\ldots ,x_{r})\det ((f_{j}^{(i)}(t_{i}))_{\substack{1\leqslant i\leqslant r\\ 0\leqslant j <   r}})
\]

(cf. FVR, I, p. 48, exerc. 11).

\begin{enumerate}
\def\labelenumi{\alph{enumi})}
\setcounter{enumi}{1}
\tightlist
\item
  Let \(r\) a positive integer and \(m_1, \ldots, m_r\) of positive integers. Let
\end{enumerate}

\[
\Delta_ {r} = \det ((x _ {i} ^ {m _ {j}}) _ {1 \leqslant i, j \leqslant r}) \in \mathbf {Z} [ x _ {1}, \dots , x _ {r} ].
\]

The polynomial \(\Delta_r\) is divisible by \(V_r\) into \(\mathbf{Z}[x_1, \ldots, x_r]\) . One denotes by \(F_r = \Delta_r / V_r\) .

\begin{enumerate}
\def\labelenumi{\alph{enumi})}
\setcounter{enumi}{2}
\tightlist
\item
  Let p be a prime number. Suppose that
\end{enumerate}

\[
0 \leqslant m _ {1} <   m _ {2} <   \dots <   m _ {r} \leqslant p - 1.
\]

Let \(\xi_{1},\ldots,\xi_{r}\) of the roots \(p\) -th roots of unity distinct. We have

\[
\mathrm{F} _ {r} (\xi_ {1}, \ldots , \xi_ {r}) \neq 0.
\]

(“Chebotarev’s Theorem”; let \(\xi\) a primitive root \(p\) th root of unity; we have \(\mathrm{F}_r(\xi_1, \ldots, \xi_r) = \mathrm{P}(\xi)\) where \(\mathrm{P} \in \mathbf{Z}[x]\) has degree \(\leqslant p - 1\) by the irreducibility of cyclotomic polynomials, the condition \(\mathrm{F}_r(\xi_1, \ldots, \xi_r) = 0\) would imply that \(\mathrm{F}_r(1, \ldots, 1)\) is divisible by \(p\) ; using \(a\) ) with \(x_i \to 1\) to prove that this is not the case.)

\begin{enumerate}
\def\labelenumi{\alph{enumi})}
\setcounter{enumi}{3}
\tightlist
\item
  Let G = Z/pZ. For every function \(f \neq 0\) from G to C, we have
\end{enumerate}

\[
\operatorname{Card} (\operatorname{Supp} (f)) + \operatorname{Card} (\operatorname{Supp} (\widehat {f})) \geqslant p + 1.
\]

\begin{enumerate}
\def\labelenumi{\alph{enumi})}
\setcounter{enumi}{4}
\tightlist
\item
  For all nonempty subsets A and B of G such that
\end{enumerate}

\[
\operatorname{Card} (\mathrm{A}) + \operatorname{Card} (\mathrm{B}) \geqslant p + 1,
\]

there exists \(f \colon \mathbf{G} \to \mathbf{C}\) whose support is A and such that the support of \(\widehat{f}\) is B. (First consider the case where \(\operatorname{Card}(\mathrm{A}) + \operatorname{Card}(\mathrm{B}) = p + 1\) .)

\begin{enumerate}
\def\labelenumi{\alph{enumi})}
\setcounter{enumi}{5}
\tightlist
\item
  Let A and B be arbitrary subsets of G. We have
\end{enumerate}

\[
\operatorname{Card} (\mathrm{A} + \mathrm{B}) \geqslant \inf (\operatorname{Card} (\mathrm{A}) + \operatorname{Card} (\mathrm{B}) - 1, p)
\]

(by the Cauchy-Davenport theorem; apply the previous question to (A, C) and to (B, D), where C and D are subsets of cardinality respectively \(p + 1 - \text{Card}(A)\) and \(p + 1 - \text{Card}(B)\) , such that \(C \cap D\) a cardinal \(\sup(\text{Card}(A) + \text{Card}(B) - p, 1)\) .)

\(g)\) Let \((a_{1},\ldots ,a_{2p - 1})\) be a family of elements of \(\mathbf{F}_p\) . There exists a subset \(\mathrm{I}\subset \{1,\dots ,2p - 1\}\) of cardinality \(p\) such that

\[
\sum_ {i \in \mathrm{I}} a _ {i} = 0.
\]

This property is false in general for a family in \(F_{p}^{2p-2}\) . (Let \(0 \leqslant b_{1} \leqslant \cdots \leqslant b_{2p-1} < p\) integers whose reductions modulo p form a permutation of \((a_{1}, \ldots, a_{2p-1})\) ; directly handle the case where there exists i such that \(a_{i} = a_{i+p-1}\) , and apply the previous question in the contrary case.)

\begin{enumerate}
\def\labelenumi{\alph{enumi})}
\setcounter{enumi}{7}
\tightlist
\item
  Let \(k \geqslant 1\) be an integer. Let \(\nu \geqslant 0\) be the integer such that \(k = p^{\nu}r\) where p does not divide r. Denote by \(\gamma\) the integer
\end{enumerate}

\[
\gamma = \left\{ \begin{array}{l l} \nu + 1 & \text { if } p \geqslant 3 \\ \nu + 2 & \text { otherwise. } \end{array} \right.
\]

If \(\ell \geqslant 4k\) is an integer, then there exist integers \((n_1, \ldots, n_\ell)\) such that \(p\) does not divide all the \(n_i\) and such that

\[
n _ {1} ^ {k} + \dots + n _ {\ell} ^ {k} \equiv 0 \mathrm{mod.} p ^ {\gamma}.
\]

\begin{enumerate}
\def\labelenumi{\arabic{enumi})}
\setcounter{enumi}{33}
\tightlist
\item
  For every compact abelian group G, we denote by \(\mu_{G}\) the normalized Haar measure on G.
\end{enumerate}

\begin{enumerate}
\def\labelenumi{\alph{enumi})}
\item
  Let G be a compact commutative group. Let \((\nu_{n})_{n\in\mathbf{N}}\) be a sequence of measures on G such that \(\nu_{n}(\mathrm{G})=1\) for all n. Then \((\nu_{n})\) converges vaguely to \(\mu_{G}\) if and only if \(\nu_{n}(\widehat{x})\) converges to 0 for every character \(\widehat{x}\neq1\) of G (“Weyl's equidistribution criterion”).
\item
  Let \(r \geqslant 1\) be an integer and \(x = (x_{1}, \ldots, x_{r}) \in \mathbf{T}^{r}\) . For \(n \geqslant 0\) , we denote by \(\nu_{n}\) be the measure on \(T^{r}\) such that
\end{enumerate}

\[
\nu_ {n} (f) = \frac {1}{n} \sum_ {k = 0} ^ {n - 1} f (k x)
\]

for \(f \in \mathcal{C}(\mathbf{T}^r)\) . The sequence \((\nu_n)\) converges vaguely to the normalized Haar measure \(\mu_{\mathrm{H}}\) , where H is the closure in \(\mathbf{T}^r\) of the set \(\{kx \mid k \in \mathbf{Z}\}\) .

\begin{enumerate}
\def\labelenumi{\alph{enumi})}
\setcounter{enumi}{2}
\tightlist
\item
  Let \(y_{i}\) of real numbers such that \(x_{i}\) is the class of \(y_{i}\) . If the family \((1, y_{1}, \ldots, y_{n})\) is Q-linearly independent, then we have H = T \(^{r}\) (“Kronecker's theorem.”)
\end{enumerate}

\(d)\) Let P be the set of prime numbers and \(\mathrm{G} = \mathbf{U}^{\mathrm{P}}\) . For every real number \(\mathrm{T} > 0\) , let \(\lambda_{\mathrm{T}}\) the measure \(\mathrm{T}^{-1}\lambda\) on \([0,\mathrm{T}]\) , where \(\lambda\) is the Lebesgue measure on \(\mathbf{R}\) . Let \(\nu_{\mathrm{T}}\) the image measure of \(\lambda_{\mathrm{T}}\) under the continuous map \(t\mapsto (p^{it})_{p\in \mathrm{P}}\) of \([0,\mathrm{T}]\) in G. When \(\mathrm{T}\to +\infty\) , the family \((\nu_{\mathrm{T}})\) converges vaguely to \(\mu_{\mathrm{G}}\) .

\begin{enumerate}
\def\labelenumi{\arabic{enumi})}
\setcounter{enumi}{34}
\tightlist
\item
  \begin{enumerate}
  \def\labelenumii{\alph{enumii})}
  \tightlist
  \item
    For every integer \(N \geqslant 1\) for every family \((a_{n})_{1 \leqslant n \leqslant N}\) of complex numbers and every integer \(H \geqslant 1\) , one has
  \end{enumerate}
\end{enumerate}

\[
\left|\sum_{n = 1}^{N}a_{n}\right|^{2}\leqslant \left(1 + \frac{N + 1}{H}\right)\sum_{|h| <   H}\left(1 - \frac{|h|}{H}\right)\sum_{\substack{1\leqslant n\leqslant N\\ 1\leqslant n + h\leqslant N}}a_{n + h}\overline{a}_{n}
\]

(“van der Corput inequality”). (Set \(a_{n}=0\) if \(n\) is not between 1 and N; write \(\sum_{n}a_{n}=\sum_{n}a_{n+h}\) for \(0\leqslant h<\mathrm{H}\) , sum over \(h\) , then apply the Cauchy-Schwarz inequality.)

We denote \(\pi\colon\mathbf{R}\to\mathbf{R}/\mathbf{Z}\) the canonical projection. A sequence \((x_{n})\) of real numbers is said to be equidistributed modulo 1 if the sequence \((\nu_{n})\) of measures on \(\mathbf{R}/\mathbf{Z}\) defined by

\[
\nu_ {n} (f) = \frac {1}{n} \sum_ {k = 1} ^ {n} f (\pi (x _ {k})), \quad f \in \mathscr {C} (\mathbf {R} / \mathbf {Z})
\]

converges vaguely to the Haar measure of total mass 1 on R/Z.

\begin{enumerate}
\def\labelenumi{\alph{enumi})}
\setcounter{enumi}{1}
\item
  If, for every integer \(h \geqslant 1\) the sequence \((x_{n+h} - x_n)_n\) is equidistributed modulo 1, then \((x_n)\) is equidistributed modulo 1. (Use the preceding exercise and the van der Corput inequality.)
\item
  Let \(d \geqslant 1\) be an integer and \(P = \alpha_{d} X^{d} + \cdots + \alpha_{0} \in R[X]\) a polynomial such that \(\alpha_{d} \notin Q\) . Then the sequence \((\mathrm{P}(n))\) is equidistributed modulo 1. (Reason by induction on \(d \geqslant 1\) .)
\end{enumerate}

\begin{enumerate}
\def\labelenumi{\arabic{enumi})}
\setcounter{enumi}{35}
\tightlist
\item
  Let \(\sigma\) a positive real number. We denote by \(\mathcal{E}_{\sigma}\) the complex vector space of entire functions \(f\) on \(\mathbf{C}\) such that, for every real number \(\varepsilon > 0\) , one has
\end{enumerate}

\[
| f (z) | = \mathrm{O} (e ^ {2 \pi (\sigma + \varepsilon) | z |}), \quad \text { as } | z | \to + \infty .
\]

\begin{enumerate}
\def\labelenumi{\alph{enumi})}
\tightlist
\item
  Suppose \(\sigma > 0\) . Let \(g \in \mathrm{L}^2([- \sigma, \sigma])\) . The function
\end{enumerate}

\[
f (z) = \int_ {- \sigma} ^ {\sigma} g (x) e ^ {- 2 i \pi x z} d x
\]

belongs to \(\mathcal{E}_{\sigma}\) The restriction of \(f\) to \(\mathbf{R}\) belongs to \(\mathrm{L}^2 (\mathbf{R})\) .

Let \(f \in \mathcal{E}_{\sigma}\) such that the restriction of \(f\) to \(\mathbf{R}\) belongs to \(\mathrm{L}^2(\mathbf{R})\) . Let \(g \in \mathrm{L}^2(\mathbf{R})\) the Fourier transform of the restriction of \(f\) to \(\mathbf{R}\) . The goal of this exercise is to prove that the support of \(g\) is contained in \([-\sigma, \sigma]\) ; by question \(a\) ), this provides a characterization of the space \(\mathcal{E}_{\sigma}\) (“Paley-Wiener theorem”).

\begin{enumerate}
\def\labelenumi{\alph{enumi})}
\setcounter{enumi}{1}
\tightlist
\item
  For every real number \(\theta\) , the function
\end{enumerate}

\[
h _ {\theta} (z) = e ^ {- i \theta} \int_ {0} ^ {+ \infty} f (r e ^ {- i \theta}) \exp (- 2 \pi z r e ^ {- i \theta}) d r
\]

is holomorphic in the half-plane \(H_{\theta}\) of C determined by the inequality

\[
\mathcal {R} (z) \cos (\theta) + \mathcal {I} (z) \sin (\theta) > \sigma .
\]

\begin{enumerate}
\def\labelenumi{\alph{enumi})}
\setcounter{enumi}{2}
\item
  If \(\theta_{1}\) and \(\theta_{2}\) are such that \(|\theta_{1} - \theta_{2}| < \pi\) , then \(h_{\theta_1}|\mathrm{H}_{\theta_1} \cap \mathrm{H}_{\theta_2} = h_{\theta_2}|\mathrm{H}_{\theta_2} \cap \mathrm{H}_{\theta_1}\) . (Apply the Cauchy formula.)
\item
  The function \(h_{0}\) (resp. \(h_{\pi}\) ) is holomorphic in the open set of \(z \in C\) such that \(\mathcal{R}(z) > 0\) (resp. \(\mathcal{R}(z) < 0\) ).
\item
  Let \(U = C - i[-\sigma, \sigma]\) . There exists a unique holomorphic function h on U such that \(h(z) = h_{0}(z)\) if \(\mathcal{R}(z) > 0\) and \(h(z) = h_{\pi}(z)\) if \(\mathcal{R}(z) < 0\) .
\item
  When \(x \rightarrow 0\) through values \textgreater{} 0, the functions
\end{enumerate}

\[
y \mapsto h _ {0} (x + i y) - h _ {\pi} (- x + i y)
\]

converge to \(g\) into \(\mathrm{L}^2 (\mathbf{R})\) .

\begin{enumerate}
\def\labelenumi{\alph{enumi})}
\setcounter{enumi}{6}
\tightlist
\item
  Conclude that g is zero almost everywhere outside \([-\sigma, \sigma]\) .
\end{enumerate}

\begin{enumerate}
\def\labelenumi{\arabic{enumi})}
\setcounter{enumi}{36}
\tightlist
\item
  We keep the notations of the preceding exercise.
\end{enumerate}

\begin{enumerate}
\def\labelenumi{\alph{enumi})}
\tightlist
\item
  Let \(f\colon\mathbf{C}\to\mathbf{C}\) be an element of \(\mathcal{E}_{1/2}\) such that the restriction of \(f\) to \(\mathbf{R}\) belongs to \(\mathrm{L}^{2}(\mathbf{R})\) . For every \(z\in\mathbf{C}-\mathbf{Z}\) , one has
\end{enumerate}

\[
f (z) = \frac {\sin (\pi z)}{\pi} \sum_ {n \in \mathbf {Z}} (- 1) ^ {n} \frac {f (n)}{z - n}.
\]

Use the Paley-Wiener theorem and the formula \(a\) ) of exercise 1.)

\begin{enumerate}
\def\labelenumi{\alph{enumi})}
\setcounter{enumi}{1}
\tightlist
\item
  Let \(f \colon \mathbf{C} \to \mathbf{C}\) be an element of \(\mathcal{E}_{1/2}\) such that \(f\) is bounded on \(\mathbf{R}\) . For every \(z \in \mathbf{C} - \mathbf{Z}\) , one has
\end{enumerate}

\[
f (z) = f (0) + \frac {\sin (\pi z)}{\pi} \Big (f ^ {\prime} (0) + \sum_ {n \in \mathbf {Z} - \{0 \}} (- 1) ^ {n} z \frac {f (n) - f (0)}{n (z - n)} \Big)
\]

(apply the previous question to the function defined by \(z \mapsto (f(z) - f(0)) / z\) for \(z \neq 0\) and which maps 0 to \(f'(0)\) ).

\begin{enumerate}
\def\labelenumi{\alph{enumi})}
\setcounter{enumi}{2}
\tightlist
\item
  Let \(f\) be an element of \(\mathcal{E}_{1/2}\) bounded on \(\mathbf{R}\) . Then \(f'\) is bounded on \(\mathbf{R}\) and we have
\end{enumerate}

\[
\sup _ {x \in \mathbf {R}} | f ^ {\prime} (x) | \leqslant 2 \sup _ {x \in \mathbf {R}} | f (x) |
\]

("Bernstein's inequality"; reduce to bounding \(|f'(1/2)|\) and apply the preceding formula).

\begin{enumerate}
\def\labelenumi{\alph{enumi})}
\setcounter{enumi}{3}
\tightlist
\item
  Let f be an element of \(E_{1/2}\) such that the restriction of f to R belongs to \(L^{2}(\mathbf{R})\) Let u and v be the 1-periodic functions on R defined by
\end{enumerate}

\[
u (t) = \widehat {f} (t) + \widehat {f} (t - 1), \quad v (t) = 2 i \pi (t \widehat {f} (t) + (t - 1) \widehat {f} (t - 1))
\]

for \(0 \leqslant t < 1\) . For almost every \(t \in [-1, 1]\) , one has

\[
\widehat {f} (t) = (1 - | t |) u (t) + (2 i \pi) ^ {- 1} \mathrm{s} (t) v (t),
\]

and the coefficients of the Fourier series of u and v are \(f(-n)\) and \(f'(-n)\) , respectively, for \(n \in Z\) .

\begin{enumerate}
\def\labelenumi{\alph{enumi})}
\setcounter{enumi}{4}
\tightlist
\item
  Let \(f \in \mathcal{E}_{1/2}\) as in the previous question. We have for all \(z \in \mathbf{C}\) the formula
\end{enumerate}

\[
f (z) = \Big (\frac {\sin (\pi z)}{\pi} \Big) ^ {2} \Big \{\sum_ {n \in \mathbf {Z}} \frac {f (n)}{(z - n) ^ {2}} + \sum_ {n \in \mathbf {Z}} \frac {f ^ {\prime} (n)}{z - n} \Big \},
\]

where the series converges symmetrically in \(\mathcal{O}(\mathbf{C})\) , that is, uniformly on compact subsets of \(\mathbf{C}\) . (Use the formulas from the question \(f\) ) of exercise 1.)

\begin{enumerate}
\def\labelenumi{\arabic{enumi})}
\setcounter{enumi}{37}
\tightlist
\item
  For every \(n \in Z\) , we denote by \(e_{n}\) the character \(x \mapsto e^{2i\pi nx}\) of T. We denote by P the subspace of \(\mathcal{C}(\mathbf{T})\) generated by the \(e_{n}\) .
\end{enumerate}

Let \(a = (a_h)_{h \in \mathbf{Z}}\) a family of complex numbers. The conjugate-harmonic family is the family \(\widetilde{a} = (-i\mathrm{s}(h)a_h)_{h \in \mathbf{Z}}\) , where \(\mathrm{s}(h)\) denotes the sign function.

For \(f \in \mathrm{P}\) , we denote by \(\widetilde{f} \in \mathrm{P}\) the function on \(\mathbf{T}\) such that \(\mathcal{F}(\widetilde{f}) = \widetilde{\mathcal{F}(f)}\) (conjugate-harmonic function of \(f\) ).

\begin{enumerate}
\def\labelenumi{\alph{enumi})}
\tightlist
\item
  Let \(f \in P\) such that \(\mathcal{F}(f)(0) = 0\) and such that f is real-valued. For every integer \(k \geqslant 1\) , there exists a constant \(C_k\) , independent of f, such that
\end{enumerate}

\[
\| \widetilde {f} \| _ {2 k} \leqslant \mathrm{C} _ {k} \| f \| _ {2 k}.
\]

(Observe that the Fourier coefficients of \(g = f + i\tilde{f}\) are supported on \(N^{*}\) , write

\[
\mathcal {R} \left(\int_ {\mathbf {T}} g (x) ^ {2 k} d x\right) = 0
\]

and expand using the binomial formula.)

\begin{enumerate}
\def\labelenumi{\alph{enumi})}
\setcounter{enumi}{1}
\item
  Let \(1 < p < +\infty\) . The map \(f \mapsto \widetilde{f}\) of P into P admits a unique continuous linear extension from \(\mathrm{L}^p(\mathbf{T})\) into \(\mathrm{L}^p(\mathbf{T})\) such that \(\mathcal{F}(\widetilde{f}) = \widetilde{\mathcal{F}(f)}\) for all \(f \in \mathrm{L}^p(\mathbf{T})\) (cf. exercise 8). (For \(p \geqslant 2\) , using the previous question and the Riesz inequality (INT, IV, §6, exerc. 18); for \(1 < p < 2\) (using duality.)
\item
  The result of the previous question does not extend to the case p = 1. (Use part (b) of Exercise 30.)
\end{enumerate}

\begin{enumerate}
\def\labelenumi{\arabic{enumi})}
\setcounter{enumi}{38}
\tightlist
\item
  Let \(\nu \in \mathcal{P}(\mathbf{R})\) a probability measure on \(\mathbf{R}\) The characteristic function of \(\nu\) is the function \(\varphi_{\nu}\) of \(\mathbf{R}\) into \(\mathbf{C}\) such that
\end{enumerate}

\[
\varphi_ {\nu} (x) = \int_ {\mathbf {R}} e ^ {i t x} d \nu (t)
\]

for all \(x \in R\) .

\begin{enumerate}
\def\labelenumi{\alph{enumi})}
\item
  For \(\nu_{1},\nu_{2}\) into \(\mathcal{P}(\mathbf{R})\) , one has \(\nu_{1} = \nu_{2}\) if and only if \(\varphi_{\nu_1} = \varphi_{\nu_2}\)
\item
  For \(\nu_{1},\nu_{2}\) into \(\mathcal{P}(\mathbf{R})\) , one has \(\nu_{1}*\nu_{2}\in \mathcal{P}(\mathbf{R})\) and \(\varphi_{\nu_1*\nu_2} = \varphi_{\nu_1}\varphi_{\nu_2}\) .
\item
  Let \(\nu \in \mathcal{P}(\mathbf{R})\) . The function \(\varphi_{\nu}\) is bounded by 1 and uniformly continuous on \(\mathbf{R}\) .
\end{enumerate}

\(d)\) Let \(\nu \in \mathcal{P}(\mathbf{R})\) For all real numbers \(a < b\) , one has

\[
\nu (] a, b [) + \frac {1}{2} \nu (\{a, b \}) = \lim _ {\mathrm{T} \rightarrow + \infty} \frac {1}{2 \mathrm{T}} \int_ {- \mathrm{T}} ^ {\mathrm{T}} \varphi_ {\nu} (x) \frac {e ^ {- i x a} - e ^ {- i x b}}{i x} d x,
\]

and in particular the limit always exists. Moreover, for every T \textgreater{} 0, we have

\[
\nu ([ - 2 \mathrm{T} ^ {- 1}, 2 \mathrm{T} ^ {- 1} ]) \leqslant \frac {1}{\mathrm{T}} \int_ {- \mathrm{T}} ^ {\mathrm{T}} \left(1 - \mathcal {R} (\varphi_ {\nu} (x))\right) d x.
\]

\begin{enumerate}
\def\labelenumi{\alph{enumi})}
\setcounter{enumi}{4}
\tightlist
\item
  Let \((\nu_{n})_{n\in\mathbf{N}}\) a sequence of probability measures on R converging narrowly (INT, IX, p.~59, def. 1) to a measure \(\nu\) ; then \(\nu\in\mathcal{P}(\mathbf{R})\) and \(\varphi_{\nu_{n}}(x)\) converges to \(\varphi_{\nu}(x)\) for all \(x\in R\) .
\end{enumerate}

Let \((\nu_{n})_{n\in\mathbf{N}}\) a sequence of probability measures on R. Suppose that for every \(x\in\mathbf{R}\) the sequence \((\varphi_{\nu_{n}}(x))_{n}\) converges to a complex number \(\varphi(x)\) , and that the function \(x\mapsto\varphi(x)\) is continuous at 0.

\begin{enumerate}
\def\labelenumi{\alph{enumi})}
\setcounter{enumi}{5}
\tightlist
\item
  Let \(\varepsilon > 0\) . There exists \(a > 0\) such that
\end{enumerate}

\[
\nu_ {n} ([ - a, a ]) \geqslant 1 - \varepsilon
\]

for every integer \(n \in N\) .

\begin{enumerate}
\def\labelenumi{\alph{enumi})}
\setcounter{enumi}{6}
\item
  The sequence \((\nu_{n})\) is relatively compact for the narrow topology.
\item
  The sequence \((\nu_{n})\) converges narrowly to a probability measure \(\nu\) whose \(\varphi\) is the characteristic function (“P. Lévy’s continuity theorem”).
\item
  There exist sequences \((\nu_{n})\) of probability measures on R, not converging tightly to a probability measure, such that \(\varphi_{\nu_{n}}(x)\) converges for all \(x \in R\) .
\end{enumerate}

\begin{enumerate}
\def\labelenumi{\arabic{enumi})}
\setcounter{enumi}{39}
\tightlist
\item
  Let \(\nu\) a probability measure on \(\mathbf{R}\) such that the identity function \(x\) of \(\mathbf{R}\) belongs to \(\mathrm{L}^2 (\nu)\) . Denote by \(m = \nu (x)\in \mathbf{R}\) and let \(\sigma \geqslant 0\) such that \(\sigma^2 = \nu ((x - m)^2)\) . We suppose \(\sigma >0\) .
\end{enumerate}

\begin{enumerate}
\def\labelenumi{\alph{enumi})}
\item
  Let \(f(t) = (t - m) / \sigma\) The image measure \(\mu = f(\nu)\) belongs to \(\mathcal{P}(\mathbf{R})\) , it satisfies \(\mu(x) = 0\) and \(\mu(x^2) = 1\) .
\item
  For \(n \geqslant 1\) , let \(\mu^{n} = \mu * \cdots * \mu\) the n-th convolution power of \(\mu\) . The sequence \((\mu^{n})_{n \geqslant 1}\) converges narrowly to the probability measure with density
\end{enumerate}

\[
\frac {1}{\sqrt {2 \pi}} e ^ {- x ^ {2} / 2}
\]

with respect to Lebesgue measure on R, called the “standard Gaussian law” (“Fundamental limit theorem of probability theory”; apply the preceding exercise).

\begin{enumerate}
\def\labelenumi{\alph{enumi})}
\setcounter{enumi}{2}
\tightlist
\item
  Let \(d \geqslant 2\) an integer. For every integer \(n \geqslant 1\) , we denote by \(s_{n}\) the map from {[}0, 1{]} into \(\{0, \ldots, d - 1\}\) which assigns to x the n-th term of the expansion of x in base d, if it is unique, or the n-th term of the limited expansion of x in base d, otherwise (TG, IV, p.~43, no.~5).
\end{enumerate}

Let \(\lambda\) the Lebesgue measure on \([0,1]\) . Denote by \(m_d = (d - 1)/2\) and \(\sigma_d = \sqrt{(d^2 - 1)/12}\) For all real numbers \(a < b\) , one has

\[
\begin{array}{r l} \lim _ {n \to + \infty} \lambda \Big (\Big \{x \in [ 0, 1 ] \mid a <   \frac {s _ {1} (x) + \cdots + s _ {n} (x) - n m _ {d}}{\sigma_ {d} \sqrt {n}} <   b \Big \} \Big) \\ & = \frac {1}{\sqrt {2 \pi}} \int_ {a} ^ {b} e ^ {- x ^ {2} / 2} d x. \end{array}
\]

\begin{enumerate}
\def\labelenumi{\arabic{enumi})}
\setcounter{enumi}{40}
\tightlist
\item
  Let \(\lambda \in \mathbf{R}_{+}^{*}\) . One denotes by \(\varpi_{\lambda}\) the discrete measure on \(\mathbf{R}\) supported on \(\mathbf{N}\) such that
\end{enumerate}

\[
\varpi_ {\lambda} (k) = e ^ {- \lambda} \frac {\lambda^ {k}}{k !}
\]

(“Poisson law with parameter λ”).

\begin{enumerate}
\def\labelenumi{\alph{enumi})}
\item
  The measure \(\varpi_{\lambda}\) is a probability measure. Compute the characteristic function of \(\varpi_{\lambda}\) ; show that the identity function x of R belongs to \(\mathrm{L}^{2}(\mathbf{R},\varpi_{\lambda})\) and compute \(\varpi(x)\) , \(\varpi(x^{2})\) .
\item
  The family of measures
\end{enumerate}

\[
\frac {\varpi_ {\lambda} - \lambda}{\sqrt {\lambda}}
\]

converges narrowly as \(\lambda\to+\infty\) to the standard Gaussian measure (cf.~preceding exercise, \(b\) ).

\begin{enumerate}
\def\labelenumi{\arabic{enumi})}
\setcounter{enumi}{41}
\tightlist
\item
  Let \(n \geqslant 2\) an integer. We denote by G the group \(S_{n}\) endowed with its Haar measure \(\mu_{n}\) of total mass 1. Let \(\varpi: S_{n} \to \mathbf{R}\) be the map which to \(\sigma\) , associates \(a + b\) , where \(a\) is the number of cycles in the decomposition of \(\sigma\) into a product of disjoint cycles (A, I, p.~60, prop. 7) and \(b\) is the number of fixed points of \(\sigma\) . One denotes by \(\nu_{n}\) the image measure \(\nu_{n} = \varpi(\mu_{n})\) on \(\mathbf{R}\) . It is a probability measure.
\end{enumerate}

\begin{enumerate}
\def\labelenumi{\alph{enumi})}
\tightlist
\item
  The characteristic function \(\varphi_{n}\) of \(\nu_{n}\) satisfies
\end{enumerate}

\[
\varphi_ {n} (t) = \prod_ {j = 1} ^ {n} \left(1 - \frac {1}{j} + \frac {e ^ {i t}}{j}\right)
\]

for all \(t \in R\) .

\begin{enumerate}
\def\labelenumi{\alph{enumi})}
\setcounter{enumi}{1}
\tightlist
\item
  Let \(\varpi_{n}\) a Poisson law with parameter \(\log(n)\) and \(\psi_{n}\) its characteristic function. When \(n \to +\infty\) , one has
\end{enumerate}

\[
\varphi_ {n} (t) = \frac {1}{\Gamma (i t)} \psi_ {n} (t) (1 + o (1))
\]

uniformly for \(t \in R\) . (Use Gauss's formula for the gamma function, cf.~FVR, VII, p.~11, formula (3).)

\begin{enumerate}
\def\labelenumi{\alph{enumi})}
\setcounter{enumi}{2}
\tightlist
\item
  Let \(f_{n}(x) = (x - \log(n))/\sqrt{\log(n)}\) for \(x \in R\) . The sequence of measures \(f_{n}(\mu_{n})\) converges narrowly to the standard Gaussian law when \(n \to +\infty\) .
\end{enumerate}

\begin{enumerate}
\def\labelenumi{\arabic{enumi})}
\setcounter{enumi}{42}
\tightlist
\item
  Let \(k \geqslant 1\) an integer. We denote \(\lambda\) the Lebesgue measure on \(R^{k}\) . Let \((\mu_{n})_{n \in \mathbf{N}}\) a sequence of probability measures on \(R^{k}\) and let \(\mu\) a probability measure on \(R^{k}\) . Suppose that the following conditions are satisfied:
\end{enumerate}

\begin{enumerate}
\def\labelenumi{(\roman{enumi})}
\item
  The Fourier transform of \(\mu\) belongs to \(\mathrm{L}^1 (\mathbf{R}^k)\) ;
\item
  There exists a sequence \((\sigma_{n})\) into \(\mathrm{GL}(k,\mathbf{R})\) such that the sequence of image measures \((\sigma_{n}(\mu_{n}))\) converges narrowly to \(\mu\) ;
\item
  For every positive real number \(\tau\) , let \(\varphi_{n,\tau}\) the characteristic function of the set of \(t \in \mathbf{R}^k\) such that \(\|^{t} \sigma_{n}(t) \| \leqslant \tau\) . We have
\end{enumerate}

\[
\lim _ {a \to + \infty} \sup _ {n \in \mathbf {N}} \int_ {\| t \| \geqslant a} | \widehat {\mu} _ {n} (^ {t} \sigma_ {n} (t)) | \varphi_ {n, \tau} (t) d t = 0.
\]

By Exercise 5, condition (i) implies that the measure \(\mu\) has a continuous density \(\psi\) with respect to Lebesgue measure. We denote by \(\alpha = \psi(0)\) . \(a)\) Let \(f \in \mathcal{C}_b(\mathbf{R}^k)\) an integrable function such that \(\widehat{f}\) has compact support. We have

\[
\lim _ {n \to + \infty} \frac {1}{| \det (\sigma_ {n}) |} \int_ {\mathbf {R} ^ {k}} f \mu_ {n} = \alpha \int_ {\mathbf {R} ^ {k}} f.
\]

(Apply the transposition formula for the Fourier transform.)

\begin{enumerate}
\def\labelenumi{\alph{enumi})}
\setcounter{enumi}{1}
\tightlist
\item
  The formula of the preceding question is valid when \(f \in \mathcal{K}(\mathbf{R}^{k})\) . For every measurable bounded subset \(B \subset R^{k}\) such that the boundary of B is negligible for Lebesgue measure, we have
\end{enumerate}

\[
\lim _ {n \to + \infty} \frac {1}{| \det (\sigma_ {n}) |} \mu_ {n} (\mathrm{B}) = \alpha \lambda (\mathrm{B}).
\]

(For every function \(f \in \mathcal{H}_{\mathbf{R}}(\mathbf{R}^k)\) , and every \(\varepsilon > 0\) , there exist integrable functions \(f_1\) and \(f_2\) whose Fourier transforms have compact support and such that \(f_1 \leqslant f \leqslant f_2\) and \(\int (f_2 - f_1) \leqslant \varepsilon\) .)

\begin{enumerate}
\def\labelenumi{\alph{enumi})}
\setcounter{enumi}{2}
\tightlist
\item
  Suppose that \(\sigma_{n}^{-1}\to 0\) in the space of real matrices of type \((k,k)\) and that there exists a continuous function \(\Phi \colon \mathbf{R}^k\to \mathbf{C}^*\) such that
\end{enumerate}

\[
\widehat {\mu} _ {n} \cdot \Phi^ {- 1} \cdot (\widehat {\mu} \circ {} ^ {t} \sigma_ {n} ^ {- 1}) ^ {- 1} \to 1
\]

uniformly on compact subsets of \(R^{k}\) . Then the above conditions are satisfied.

\begin{enumerate}
\def\labelenumi{\alph{enumi})}
\setcounter{enumi}{3}
\tightlist
\item
  Take k = 1. Let \(\mu\) be a probability measure on R with base Lebesgue measure such that the identity function x of R belongs to \(L^{2}(\nu)\) and such that \(\mu(x) = 0\) . Let \(\sigma \geqslant 0\) such that \(\sigma^{2} = \mu(x^{2})\) . We suppose \(\sigma > 0\) . For every integer \(n \geqslant 1\) , let \(\mu^{n}\) the n-th convolution power of \(\mu\) . For every bounded measurable subset \(B \subset R\) whose boundary has measure zero with respect to Lebesgue measure, we have
\end{enumerate}

\[
\lim _ {n \rightarrow + \infty} \sigma \sqrt {n} \mu^ {n} (\mathrm{B}) = \frac {1}{2 \sqrt {\pi}} \lambda (\mathrm{B})
\]

("fundamental local limit theorem").

\begin{enumerate}
\def\labelenumi{\arabic{enumi})}
\setcounter{enumi}{43}
\tightlist
\item
  Let I = {[}a, b{]} be a compact interval of R. Let f be a continuous function from I into \(R^{2}\) which is injective on the interior of I, and which is \(C^{1}\) piecewise (that is, there exists a finite partition of I into intervals J such that f\textbar J is differentiable, and \(f(f|\mathrm{J})'\) is continuous). We say that \(f\) is a parametrized curve in the plane. We say that
\end{enumerate}

\[
\ell (f) = \int_ {\mathrm{I}} | f ^ {\prime} (t) | d t
\]

(where \(f'(t)\) is identified with a complex number) is the length of the parametrized curve \(f\) . If moreover \(f(a) = f(b)\) , we say that \(f\) defines a Jordan curve. There then exists a unique bounded connected component C of \(\mathbf{R}^2 - f(\mathrm{I})\) (TA, to appear). The Lebesgue measure of C is called the area enclosed by the Jordan curve \(f\) , and is denoted \(\mathrm{A}(f)\) .

\begin{enumerate}
\def\labelenumi{\alph{enumi})}
\tightlist
\item
    Suppose that \(f\) is a Jordan curve. For every \(t \in I\) , the restriction \(f_t\) of \(f\) to \([a, t]\) is a parametrized curve and the function \(s: t \mapsto \ell(f_t)\) is a strictly increasing bijection from I into \(J = [0, \ell(f)]\) . The function \(g = f \circ s^{-1}: J \to \mathbf{R}^2\) is a Jordan curve, and for every \(t \in J\) , one has \(\ell(g_t) = t\) . b) Let \(f = (f_1, f_2)\) be a Jordan curve. We have
\end{enumerate}

\[
\mathrm{A} (f) = \frac {1}{2} \int_ {\mathrm{I}} (f _ {1} f _ {2} ^ {\prime} - f _ {2} f _ {1} ^ {\prime}),
\]

where the integral is with respect to Lebesgue measure on I. (Apply Stokes' formula.) Moreover, we have \(A(g) = A(f)\) , where g is the function of question a).

\begin{enumerate}
\def\labelenumi{\alph{enumi})}
\setcounter{enumi}{2}
\tightlist
\item
  We have \(4\pi\mathrm{A}(f)\leqslant\ell(f)^{2}\) ("isoperimetric inequality in the plane"). (Reduce to the case where \(\ell(f)=1\) ; consider the Fourier series of the real and imaginary parts of g and show that \(|g'(t)|^{2}=1\) for \(t\in J\) .)
\end{enumerate}

\(\text{d)}\quad\text{If }4\pi\mathrm{A}(f)=\ell(f)^{2},\text{ then the image of }f\text{ is a circle in }\mathbf{R}^{2}.\)

\begin{enumerate}
\def\labelenumi{\arabic{enumi})}
\setcounter{enumi}{44}
\tightlist
\item
  For every real number \(h > 0\) , let \(\Delta_h\) the continuous function on \(\mathbf{R}\) such that \(\Delta_h(t) = 0\) if \(|t| > h\) and \(\Delta_h(t) = 1 - |t| / h\) for \(t \in [-h, h]\) .
\end{enumerate}

\begin{enumerate}
\def\labelenumi{\alph{enumi})}
\tightlist
\item
  The family in \(\mathcal{C}([0,1])\) consisting of the functions \(f_{0}=1\) , \(f_{1}:t\mapsto t\) , and the functions
\end{enumerate}

\[
f _ {k, l} \colon t \mapsto \Delta_ {2 ^ {- k}} \left(t - \frac {2 l + 1}{2 ^ {k}}\right)
\]

for \(k \geqslant 1\) and \(0 \leqslant l < 2^{k-1}\) is a Banach basis of the Banach space \(\mathcal{C}([0,1])\) (cf.~EVT, IV, p.~70, exerc. 14). More precisely, for every \(f \in \mathcal{C}([0,1])\) , one has

\[
\begin{array}{l} f = f (0) f _ {0} + (f (1) - f (0)) f _ {1} + \\ \sum_ {k \geqslant 1} \sum_ {0 \leqslant l <   2 ^ {k - 1}} \left(f \left(\frac {2 l + 1}{2 ^ {k}}\right) - \frac {1}{2} \left(f \left(\frac {l}{2 ^ {k - 1}}\right) + f \left(\frac {l + 1}{2 ^ {k - 1}}\right)\right)\right) f _ {k, l} \end{array}
\]

into \(\mathcal{C}([0,1])\) .

\begin{enumerate}
\def\labelenumi{\alph{enumi})}
\setcounter{enumi}{1}
\tightlist
\item
  The family of functions \(t \mapsto e^{2i\pi ht}\) for \(h \in Z\) is not a Schauder basis of the space of functions \(f \in \mathcal{C}([0,1])\) such that \(f(0) = f(1)\) .
\end{enumerate}

\begin{enumerate}
\def\labelenumi{\arabic{enumi})}
\setcounter{enumi}{45}
\tightlist
\item
  For every integer \(k \geqslant 0\) , let \(\mathrm{D}_k \subset [0,1]\) the set of numbers \(j/2^k\) for \(0 \leqslant j \leqslant 2^k\) integer. Let D be the union of the sets \(\mathrm{D}_k\) for \(k \geqslant 0\) . For every \(k \geqslant 1\) , and every \(t \in \mathrm{D}_{k+1} - \mathrm{D}_k\) , let \(t_+\) (resp. \(t_-\) ) the smallest element of \(\mathrm{D}_{k+1}\) such that \(t < t_+\) (resp. the greatest element of \(\mathrm{D}_{k+1}\) such that \(t_- < t\) ); \(t_+\) and \(t_-\) belong to \(\mathrm{D}_k\) .
\end{enumerate}

\begin{enumerate}
\def\labelenumi{\alph{enumi})}
\tightlist
\item
  Let \(\sigma\colon\mathrm{D}\to[0,1]\) a strictly increasing function such that \(\sigma(0)=0\) and \(\sigma(1)=1\) whose image is dense in \([0,1]\) . Then there exists a unique strictly increasing homeomorphism from \([0,1]\) into itself that extends \(\sigma\) .
\end{enumerate}

Let \(f\colon [0,1]\to \mathbf{R}\) a continuous function. We call a level sketch \(k\) of \(f\) a strictly increasing mapping \(\sigma \colon \mathrm{D}_k\to [0,1]\) such that \(\sigma (0) = 0\) , \(\sigma (1) = 1\) and such that for every integer \(j\) such that \(0\leqslant j < 2^k\) , one of the following conditions holds:

\begin{enumerate}
\def\labelenumi{(\roman{enumi})}
\item
  \(f(\sigma(j/2^{k})) \neq f(\sigma((j+1)/2^{k}))\) ;
\item
  there exists an open interval I contained in \(\sigma(j/2^{k}), \sigma(j+1)/2^{k}\) such that, for every \(t \in I\) , one has
\end{enumerate}

\[
f (t) = f (\sigma (j / 2 ^ {k})).
\]

We denote \(t_{\sigma}\) the smallest strictly positive element of \(D_{k}\) such that

\[
\inf _ {t ^ {\prime} <   t _ {\sigma}} | \sigma (t _ {\sigma}) - \sigma (t ^ {\prime}) | = \sup _ {t \in \mathrm{D} _ {k}} \inf _ {t ^ {\prime} <   t} | \sigma (t) - \sigma (t ^ {\prime}) |.
\]

One says that a sketch \(\tau\) of level \(k + 1\) is a good extension of a draft \(\sigma\) of level \(k\) if \(\tau\) extends \(\sigma\) and if, for all \(t \in \mathrm{D}_{k+1} - \mathrm{D}_k\) , one has

\[
f (\tau (t)) = \frac {1}{2} (f (\sigma (t _ {-})) + f (\sigma (t _ {+})))
\]

if \(t_{-}\neq t_{\sigma}\)

\begin{enumerate}
\def\labelenumi{\alph{enumi})}
\setcounter{enumi}{1}
\tightlist
\item
  There exists an increasing homeomorphism \(\sigma\) of \([0,1]\) of itself into itself such that
\end{enumerate}

\begin{enumerate}
\def\labelenumi{(\roman{enumi})}
\item
  For every integer \(k \geqslant 1\) , the restriction \(\sigma_{k}\) of \(\sigma\) to \(\mathrm{D}_{k}\) is a draft of level \(k\) ;
\item
  For every integer \(k \geqslant 1\) , the map \(\sigma_{k+1}\) is a good extension of \(\sigma_{k}\) .
\end{enumerate}

\begin{enumerate}
\def\labelenumi{\alph{enumi})}
\setcounter{enumi}{2}
\tightlist
\item
  Suppose that \(f(0) = f(1)\) . Denote by \(\sigma\) the increasing homeomorphism of {[}0,1{]} defined by a mapping \(\sigma: D \to [0,1]\) satisfying the conditions of the previous question. There exists a constant \(C \geqslant 0\) such that the continuous function \(f \circ \sigma\) identified with a continuous function on T, satisfies
\end{enumerate}

\[
| \widehat {(f \circ \sigma)} (n) | \leqslant \frac {\mathrm{C}}{| n |}
\]

for all \(n \in Z - \{0\}\) ("Sahakyan's Theorem"; use Exercise 2 of II, p.~262 and the preceding exercise.)

\begin{enumerate}
\def\labelenumi{\alph{enumi})}
\setcounter{enumi}{3}
\item
  There exists an increasing homeomorphism \(\sigma\) of {[}0, 1{]} into itself such that the Fourier series of \(f \circ \sigma\) converges symmetrically in \(\mathcal{C}(\mathbf{T})\) (“Bohr-Pál Theorem”).
\item
  Does Sahakyan's theorem hold for complex-valued continuous functions?
\item
  Is there a function \(f: R/Z \to C\) whose image has nonempty interior and which satisfies the condition \(\sup_{h \in Z - \{0\}} |h||\widehat{f}(h)| < +\infty\) ?
\end{enumerate}

\begin{enumerate}
\def\labelenumi{\arabic{enumi})}
\setcounter{enumi}{46}
\tightlist
\item
  Let X be a locally compact topological space and \(\mu\) a probability measure on X. Let \((f_n)_{n\in \mathbf{N}^*}\) an orthonormal family in \(\mathrm{L}^2 (\mathrm{X},\mu)\) . Let \((a_{n})_{n\in \mathbf{N}^{*}}\) a sequence of complex numbers such that the family \((a_{n}\log (n))_{n\in \mathbf{N}^{*}}\) belongs to \(\ell^2 (\mathbf{N}^*)\) .
\end{enumerate}

For every finite set of integers I, we denote by

\[
s _ {\mathrm{I}} = \sum_ {n \in \mathrm{I}} a _ {n} f _ {n},
\]

and we also write \(s_{n} = s_{\{1,\ldots,n\}}\) .

\begin{enumerate}
\def\labelenumi{\alph{enumi})}
\tightlist
\item
  Let \(N \geqslant 1\) There exists a family I of nonempty intervals of \(\{1, \ldots, N\}\) and a constant \(C \geqslant 0\) independent of N, satisfying the following conditions:
\end{enumerate}

\begin{enumerate}
\def\labelenumi{(\roman{enumi})}
\item
  Every interval of \(\{1, \ldots, N\}\) of the form \(I = \{1, \ldots, M\}\) is a disjoint union of at most C log(M + 1) intervals in \(\mathcal{I}\) ;
\item
  Every integer \(n \in \{1, \ldots, N\}\) belongs to at most \(C \log(N + 1)\) intervals in I.
\end{enumerate}

\begin{enumerate}
\def\labelenumi{\alph{enumi})}
\setcounter{enumi}{1}
\tightlist
\item
  There exists a constant C \textgreater{} 0 such that for all \(N \in N^{*}\) , one has
\end{enumerate}

\[
\int_ {X} \left(\sup _ {1 \leqslant n \leqslant N} | s _ {n} | ^ {2}\right) d \mu \leqslant C \sum_ {n = 1} ^ {N} | a _ {n} | ^ {2}.
\]

(Let \(\tau(x)\) the smallest integer \(\leqslant\) N such that \(\sup_{1\leqslant n\leqslant N}|s_n(x)| = |s_{\tau(x)}(x)|\) , and let \(S(x) = s_{\tau(x)}(x)\) . Let \(\mathcal{I}\) a family of intervals as in the previous question; write

\[
\mathrm{S} = \sum_ {\mathrm{I} \in \mathcal {J}} s _ {\mathrm{I}}
\]

where J is a disjoint family of intervals in I, and then apply the Cauchy–Schwarz inequality.)

\begin{enumerate}
\def\labelenumi{\alph{enumi})}
\setcounter{enumi}{2}
\tightlist
\item
  For \(j\in \mathbf{N}^*\) , set
\end{enumerate}

\[
\mathrm{S} _ {j} = \sum_ {2 ^ {j} <   n \leqslant 2 ^ {j + 1}} a _ {n} f _ {n}.
\]

The series with general term \((j+1)|S_{j}|^{2}\) converges \(\mu\) -almost everywhere in X.

\(d)\) The sequence \((s_{2^j}(x))_{j\in \mathbf{N}^*}\) is a Cauchy sequence for \(\mu\) -almost all \(x\in X\) . Denote by \(f(x)\) its limit, defined \(\mu\) -almost everywhere.

\begin{enumerate}
\def\labelenumi{\alph{enumi})}
\setcounter{enumi}{4}
\tightlist
\item
  The sequence \((s_n(x))_{n\in \mathbf{N}^*}\) converges to \(f(x)\) for \(\mu\) -almost all \(x\in X\) .
\end{enumerate}

\(f)\) Let \(f\in \mathrm{L}^1 (\mathbf{T})\) such that

\[
\sum_ {h \in \mathbf {Z}} \log (| h | + 1) ^ {2} | \widehat {f} (h) | ^ {2} <   + \infty .
\]

For almost all \(x\) , the Fourier series of \(f\) to \(x\) converge symmetrically to \(f(x)\) ("Rademacher-Menshov theorem").

\begin{enumerate}
\def\labelenumi{\arabic{enumi})}
\setcounter{enumi}{47}
\tightlist
\item
  \begin{enumerate}
  \def\labelenumii{\alph{enumii})}
  \tightlist
  \item
    Let \(p \in [1, +\infty[\) A subset X of \(L^p(\mathbf{R})\) is relatively compact if and only if it is bounded and satisfies the conditions
  \end{enumerate}
\end{enumerate}

\[
\lim _ {y \to 0} \sup _ {f \in X} \| \boldsymbol {\gamma} (y) f - f \| _ {p} = 0
\]

\[
\lim _ {\mathrm{R} \to + \infty} \sup _ {f \in \mathrm{X}} \| (1 - \varphi_ {\mathrm{R}}) f \| _ {p} = 0,
\]

where \(\varphi_{R}\) is the characteristic function of the interval \([-R,R]\) (cf.~INT, VIII, p.~205, exercise 26).

\begin{enumerate}
\def\labelenumi{\alph{enumi})}
\setcounter{enumi}{1}
\tightlist
\item
  Let X be a bounded subset of \(\mathrm{L}^{2}(\mathbf{R})\) . Prove that X is precompact if and only if
\end{enumerate}

\[
\lim _ {\mathrm{R} \to + \infty} \sup _ {f \in \mathrm{X}} \int_ {\mathbf {R}} (1 - \varphi_ {\mathrm{R}}) (| f | ^ {2} + | \widehat {f} | ^ {2}) d x = 0.
\]

\begin{enumerate}
\def\labelenumi{\alph{enumi})}
\setcounter{enumi}{2}
\tightlist
\item
  Let \((f_{n})\) a sequence in \(\mathrm{L}^{2}(\mathbf{R})\) orthonormal in \(\mathrm{L}^{2}(\mathbf{R})\) such that \(\sup_{n}|f_{n}|\) belongs to \(\mathrm{L}^{2}(\mathbf{R})\) . Then \(\sup_{n}|\widehat{f_{n}}|\) does not belong to \(\mathrm{L}^{2}(\mathbf{R})\) .
\end{enumerate}

\begin{enumerate}
\def\labelenumi{\arabic{enumi})}
\setcounter{enumi}{48}
\tightlist
\item
  Let \(k\) a finite field of characteristic \(p\) and cardinality \(q\) . We endow \(k\) of counting measure and we denote \(\widehat{f}\) the inverse Fourier transform of a function \(f: k \to \mathbf{C}\) .
\end{enumerate}

\begin{enumerate}
\def\labelenumi{\alph{enumi})}
\item
  Let \(\chi\) a non-trivial character of \(k^{*}\) . We extend \(\chi\) to \(k\) by setting \(\chi(0) = 0\) Prove that \(|\widehat{\chi}(\psi)| = \sqrt{q}\) for all \(\psi \in \widehat{k}\) such that \(\psi \neq 1\) . (“Gauss sums.”)
\item
  Let \(d \mid q - 1\) an integer. Prove that for every character \(\psi \neq 1\) of \(k\) , one has
\end{enumerate}

\[
\left| \sum_ {x \in k} \psi (x ^ {d}) \right| \leqslant (d - 1) \sqrt {q},
\]

with equality if d = 2.

\begin{enumerate}
\def\labelenumi{\alph{enumi})}
\setcounter{enumi}{2}
\tightlist
\item
  Let \(\chi_{1}\) and \(\chi_{2}\) of the non-trivial characters of \(k^{*}\) , extended by 0 to k as in question a). We set
\end{enumerate}

\[
\mathrm{J} \left(\chi_ {1}, \chi_ {2}\right) = \sum_ {x \in k} \chi_ {1} (x) \chi_ {2} (1 - x)
\]

(“Jacobi sums”). If \(\chi_1\chi_2\neq 1\) , then for every character \(\psi \neq 1\) of \(k\) , one has

\[
\mathrm{J} (\chi_ {1}, \chi_ {2}) = \frac {\widehat {\chi} _ {1} (\psi) \widehat {\chi} _ {2} (\psi)}{\widehat {\chi_ {1} \chi_ {2}} (\psi)}, \qquad | \mathrm{J} (\chi_ {1}, \chi_ {2}) | = \sqrt {q}.
\]

\begin{enumerate}
\def\labelenumi{\alph{enumi})}
\setcounter{enumi}{3}
\item
  Suppose that p is a prime congruent to 1 modulo 4. There exist integers a and b such that \(a^{2} + b^{2} = p\) (“Fermat's two-square theorem”). (Consider a Jacobi sum for characters of suitably chosen order of \(F_{p}^{*}\) .)
\item
  Suppose that p is odd. Let \(\lambda \in \widehat{k}^{*}\) be the unique character of order 2 (“Legendre character”). Let \(\chi\) be a non-trivial character of \(k^{*}\) and set \(\eta = \chi^{2}\) . Prove that for every character \(\psi \neq 1\) of k, we have
\end{enumerate}

\[
\widehat {\eta} (\psi) \widehat {\lambda} (\psi) = \chi (4) \widehat {\chi} (\psi) \widehat {\chi \lambda} (\psi)
\]

(“Hasse–Davenport formula”).

\begin{enumerate}
\def\labelenumi{\alph{enumi})}
\setcounter{enumi}{5}
\tightlist
\item
  Suppose that p is odd. Let \(N_{p}\) the number of \((x,y)\in\mathbf{F}_{p}\times\mathbf{F}_{p}\) such that \(y^{2}=x^{4}+1\) . We have
\end{enumerate}

\[
\left| \mathrm{N} _ {p} - p \right| \leqslant 2 \sqrt {p}.
\]

(When \(p\) is congruent to 3 modulo 4, show that \(\mathrm{N}_p = p\) ; otherwise, verify that \(\mathrm{N}_p - p = -\sum_x \lambda(x^4 + 1)\) , where \(\lambda\) is the Legendre character modulo \(p\) , and express this sum in terms of Jacobi sums.)

\begin{enumerate}
\def\labelenumi{\alph{enumi})}
\setcounter{enumi}{6}
\tightlist
\item
  Compare these formulas with those satisfied by the gamma function (FVR, VII).
\end{enumerate}

\begin{enumerate}
\def\labelenumi{\arabic{enumi})}
\setcounter{enumi}{49}
\tightlist
\item
  Let \(k\) a finite field of characteristic \(p \neq 2\) and cardinality \(q\) . We endow \(k\) of counting measure and we denote \(\widehat{f}\) the inverse Fourier transform of a function \(f: k \to \mathbf{C}\) . We endow \(k^*\) of the counting measure and fix a nontrivial character \(\psi\) of \(k\) . One denotes by \(\lambda\) the unique character of order 2 of \(k^*\) , extended to \(k\) by setting \(\lambda(0) = 0\) .
\end{enumerate}

For every \(a \in k^{*}\) , set

\[
\mathrm{S} (a) = \sum_ {x \in k ^ {*}} \lambda (x) \psi (a x + x ^ {- 1})
\]

(“Salié sums”). We denote by \(\mathcal{F}(\mathrm{S})\) the Fourier transform of the function S on \(k^{*}\) .

\begin{enumerate}
\def\labelenumi{\alph{enumi})}
\tightlist
\item
  Let \(\chi \in \widehat{k}^*\) and set \(\eta = \chi^2\) . We have
\end{enumerate}

\[
\mathcal {F} (\mathrm{S}) (\chi) = \widehat {\overline {{\chi}}} (\psi) \widehat {\overline {{\chi}}} \lambda (\psi) = \chi (4) \widehat {\lambda} (\psi) \widehat {\overline {{\eta}}} (\psi)
\]

(use the Hasse-Davenport formula).

\begin{enumerate}
\def\labelenumi{\alph{enumi})}
\setcounter{enumi}{1}
\tightlist
\item
  Deduce that
\end{enumerate}

\[
\operatorname{S}(a) = \lambda (a)\widehat{\lambda} (\psi)\sum_{\substack{y\in k^{*}\\ y^{2} = 4a}}\psi (y),\quad |\operatorname{S}(a)|\leqslant 2\sqrt{q}.
\]

\begin{enumerate}
\def\labelenumi{\alph{enumi})}
\setcounter{enumi}{2}
\tightlist
\item
  Analogy with Bessel functions \(J_{k+1/2}\) (exercise 20) for \(k \in N\) .
\end{enumerate}

\begin{enumerate}
\def\labelenumi{\arabic{enumi})}
\setcounter{enumi}{50}
\tightlist
\item
  We identify \(\mathbf{T}\) and the interval {[}0,1{[}. For \(n \in \mathbf{Z}\) , let \(e_n\) the unitary character of \(\mathbf{T}\) defined by \(t \mapsto \exp(2i\pi nt)\) . For every measure \(\mu \in \mathcal{M}^1(\mathbf{T})\) and every integer \(m \geqslant 1\) , we denote by
\end{enumerate}

\[
\mathrm{S} _ {m} (\mu) = \sum_ {h = - m} ^ {m} \widehat {\mu} (h) e _ {h} \in \mathcal {C} (\mathbf {T})
\]

the m-th symmetric partial sum of the Fourier series of \(\mu\) .

A finite family \((y_{1},\ldots,y_{k})\) of real numbers is said to be T-independent if \((1,y_{1},\ldots,y_{k})\) is Q-linearly independent. This amounts to saying that the equation

\[
\sum_ {j = 1} ^ {k} m _ {j} y _ {j} = 0
\]

into \(\mathbf{T}\) , with \(m_j \in \mathbf{Z}\) , has as its unique solution the zero family. For \(t \in \mathbf{R}\) , we denote by \(t - (y_1, \ldots, y_k)\) the family \((t - y_j)_j\) .

\begin{enumerate}
\def\labelenumi{\alph{enumi})}
\tightlist
\item
  Let \(n \geqslant 10^{6}\) an integer. There exist real numbers \(x_{1}, \ldots, x_{n}\) into {[}0,1{]} such that
\end{enumerate}

\[
\left| x _ {j} - j / n \right| \leqslant 1 0 ^ {- 4} n ^ {- 1}
\]

and such that the family \(\mathrm{E}_{n}=(x_{1},\ldots,x_{n})\) is T-independent.

We denote

\[
\mu_ {n} = \frac {1}{n} \sum_ {j = 1} ^ {n} \varepsilon_ {x _ {j}}.
\]

This is a probability measure on T.

\begin{enumerate}
\def\labelenumi{\alph{enumi})}
\setcounter{enumi}{1}
\tightlist
\item
  Let \(t\in \mathbf{R}\) . We have
\end{enumerate}

\[
\frac {1}{n} \sum_ {j = 1} ^ {n} \frac {1}{| \sin (\pi (t - x _ {j})) |} \geqslant \frac {1}{5} \log n.
\]

\begin{enumerate}
\def\labelenumi{\alph{enumi})}
\setcounter{enumi}{2}
\tightlist
\item
  Let \(t \in R\) such that the family \(t - E_{n}\) is T-independent. There exists \(m \in N^{*}\) such that
\end{enumerate}

\[
| \mathrm{S} _ {m} (\mu_ {n}) (t) | \geqslant \frac {2}{1 5} \log n.
\]

(apply Kronecker's theorem, cf.~exercise 34).

\(d)\) Let \(t \in \mathbf{R}\) such that the family \(t - \mathrm{E}_n\) is not \(\mathbf{T}\) -independent and let \(m_1, \ldots, m_n\) be integers, not all zero, such that

\[
\sum_ {j = 1} ^ {n} m _ {j} (t - x _ {j}) = 0
\]

in T. If two of the integers \(m_{j}\) are nonzero, there exists \(m \geqslant 0\) such that

\[
| \mathrm{S} _ {m} (\mu_ {n}) (t) | \geqslant \frac {1}{1 0} \log n
\]

(let j be such that \(m_{j} \neq 0\) and \(|x_{j} - t|\) is maximal; note that the family \(t - (x_{i})_{i \neq j}\) is T-independent and apply Kronecker's theorem).

\begin{enumerate}
\def\labelenumi{\alph{enumi})}
\setcounter{enumi}{4}
\tightlist
\item
  We keep the hypothesis from the previous question, but we assume that only one of the coefficients \(m_{j}\) is nonzero, and that \(t \notin E_{n}\) . There exists \(m \geqslant 0\) such that
\end{enumerate}

\[
| \mathrm{S} _ {m} (\mu_ {n}) (t) | \geqslant \frac {1}{1 0} \log n
\]

(observe that the family \(m_{j}(t - x_{i})_{i \neq j}\) is T-independent and apply Kronecker's theorem).

\begin{enumerate}
\def\labelenumi{\alph{enumi})}
\setcounter{enumi}{5}
\tightlist
\item
  Suppose that \(t \in E_{n}\) . There exists \(m \geqslant 0\) such that
\end{enumerate}

\[
| \mathrm{S} _ {m} (\mu_ {n}) (t) | \geqslant \frac {1}{1 0} \log n
\]

(the family \((x_{i})_{x_{i}\neq t}\) is T-independent).

\begin{enumerate}
\def\labelenumi{\alph{enumi})}
\setcounter{enumi}{6}
\tightlist
\item
  There exists an integer \(\mathrm{M} \geqslant 0\) such that, for every \(t \in \mathbf{T}\) , one has
\end{enumerate}

\[
\sup _ {0 \leqslant m \leqslant \mathrm{M}} | \mathrm{S} _ {m} (\mu_ {n}) (t) | \geqslant \frac {1}{2 0} \log n.
\]

\begin{enumerate}
\def\labelenumi{\alph{enumi})}
\setcounter{enumi}{7}
\tightlist
\item
  There exists a function \(f_{n}\in\mathcal{C}^{\infty}(\mathbf{T})\) and an integer \(\mathrm{M}\geqslant0\) such that
\end{enumerate}

\[
\begin{array}{c} \operatorname{Supp} (f _ {n}) \subset \bigcup_ {j = 1} ^ {n} \Bigl [ \frac {j}{n} - \frac {2}{1 0 ^ {4} n}, \frac {j}{n} + \frac {2}{1 0 ^ {4} n} \Bigr ] \\ \int_ {j / n - 2 / (1 0 ^ {4} n)} ^ {j / n + 2 / (1 0 ^ {4} n)} f _ {n} = \frac {1}{n}, \qquad 1 \leqslant j \leqslant n \\ \sup _ {0 \leqslant m \leqslant M} | S _ {m} (f _ {n}) (t) | \geqslant \frac {1}{4 0} \log n, \qquad \text { for all } t \in \mathbf {T}. \end{array}
\]

\begin{enumerate}
\def\labelenumi{\roman{enumi})}
\tightlist
\item
  There exists \(M \geqslant 1\) and a finite family \((a_{n})_{|n| \leqslant M}\) of complex numbers such that the function
\end{enumerate}

\[
g _ {\mathrm{M}} (t) = \sum_ {k = - \mathrm{M}} ^ {\mathrm{M}} a _ {k} e _ {k} n
\]

satisfies

\[
\int_ {\mathbf {T}} | g _ {\mathrm{M}} (t) | d t \leqslant 2, \quad \inf _ {t \in \mathbf {T}} \sup _ {0 \leqslant m \leqslant \mathrm{M}} | \mathrm{S} _ {m} (g _ {\mathrm{M}}) (t) | \geqslant \frac {1}{4 0} \log n.
\]

\begin{enumerate}
\def\labelenumi{\alph{enumi})}
\setcounter{enumi}{9}
\tightlist
\item
  There exists \(\varphi\in\mathrm{L}^{1}(\mathbf{T})\) such that \(\widehat{\varphi}(h)=0\) if \(h<0\) and such that the series
\end{enumerate}

\[
\sum_ {h \geqslant 1} \widehat {\varphi} (h) e _ {h} (t)
\]

diverges for every \(t \in \mathbf{T}\) (“Kolmogorov’s theorem”).

\[
g _ {k} = \sum_ {n = - \mathrm{N} _ {k}} ^ {\mathrm{N} _ {k}} a _ {k, n} e _ {n}
\]

(There exists a sequence \(\mathbf{N}_k\) of integers and a sequence of functions

such that \(\| g_k\| _1\leqslant 2\) and

\[
\inf _ {t \in \mathbf {T}} \sup _ {0 \leqslant m \leqslant \mathrm{N} _ {k}} | \mathrm{S} _ {m} (g _ {k}) (t) | \geqslant 2 ^ {2 k}.
\]

Let \(\widetilde{g}_k = e_{\mathrm{M}_k}g_k\) , where \(\mathrm{M}_k\) is a sufficiently large integer; set \(\varphi = \sum_{k\geqslant 1}2^{-k}\widetilde{g}_k.\)

\begin{enumerate}
\def\labelenumi{\arabic{enumi})}
\setcounter{enumi}{51}
\tightlist
\item
  Let H be a countable discrete group, not necessarily commutative. a) The group H is amenable (EVT, IV, p.~73, exercise 4) if and only if there exists an increasing sequence \((\mathrm{F_N})_{\mathrm{N}}\) of finite subsets of H such that, for every \(x \in \mathrm{H}\) , one has
\end{enumerate}

\[
\lim _ {N \to + \infty} \frac {1}{\operatorname{Card} (F _ {N})} \operatorname{Card} ((F _ {N} - x F _ {N}) \cup (x F _ {N} - F _ {N})) = 0
\]

(“Følner sequence”). If H is commutative, then it is amenable.

\begin{enumerate}
\def\labelenumi{\alph{enumi})}
\setcounter{enumi}{1}
\tightlist
\item
  Let G be a compact commutative group. Fix a sequence \((\mathrm{F}_{\mathrm{N}})_{\mathrm{N}\geqslant1}\) of finite subsets of the discrete group \(\widehat{G}\) as in a) and denote \(\varphi_{N}\) be the characteristic function of \(F_{N}\) . Set
\end{enumerate}

\[
\psi_ {\mathrm{N}} = \frac {1}{\operatorname{Card} \left(\mathrm{F} _ {\mathrm{N}}\right)} \varphi_ {\mathrm{N}} * \widetilde {\varphi} _ {\mathrm{N}}.
\]

The function \(\psi_{N}\) is finitely supported, it satisfies \(0 \leqslant \psi_{N} \leqslant 1\) , and

\[
\lim _ {\mathrm{N} \to + \infty} \psi_ {\mathrm{N}} (\chi) = 1
\]

for all \(\chi\in\widehat{\mathbf{G}}\) .

\begin{enumerate}
\def\labelenumi{\alph{enumi})}
\setcounter{enumi}{2}
\tightlist
\item
  Let \(\mu_{\mathrm{N}}\) the measure on G with density the continuous map
\end{enumerate}

\[
\sum_ {\chi \in \widehat {\mathrm{G}}} \psi_ {\mathrm{N}} (\chi) \chi .
\]

The sequence of measures \((\mu_{\mathrm{N}})_{\mathrm{N}\geqslant1}\) converges to \(\varepsilon_{e}\) in the space \(\mathcal{M}^{1}(\mathrm{G})\) endowed with the topology of compact convergence in \(\mathcal{C}(\mathrm{G})\) (Apply Lemma 4 of INT, VIII, §2, no. 7.)

\begin{enumerate}
\def\labelenumi{\alph{enumi})}
\setcounter{enumi}{3}
\tightlist
\item
  Let \(f\in \mathcal{C}(\mathrm{G})\) . We then have
\end{enumerate}

\[
f = \lim _ {N \to \infty} \sum_ {\chi \in \widehat {G}} \psi_ {N} (\chi) \mathcal {F} (f) (\chi) \chi
\]

into \(\mathcal{C}(\mathrm{G})\) .

\begin{enumerate}
\def\labelenumi{\alph{enumi})}
\setcounter{enumi}{4}
\tightlist
\item
  Retrieve Fejér's theorem (prop. 23 of II, p.~242) as a special case of this statement.
\end{enumerate}

\begin{enumerate}
\def\labelenumi{\arabic{enumi})}
\setcounter{enumi}{52}
\tightlist
\item
  Let \(f\in \mathcal{L}^1 (\mathbf{T})\)
\end{enumerate}

\begin{enumerate}
\def\labelenumi{\alph{enumi})}
\item
  Let \(k \geqslant 1\) an integer and suppose that \(f \in \mathcal{C}^{k}(\mathbf{T})\) . The Fourier series of f converges to f in \(\mathcal{C}^{k-1}(\mathbf{T})\) . If \(k \geqslant 2\) , the Fourier series of f converges absolutely to f in \(\mathcal{C}(\mathbf{T})\) .
\item
  For every integer \(N \geqslant 1\) , let \(f_{N}\) the function on T such that
\end{enumerate}

\[
f_{\mathrm{N}}(t) = \sum_{\substack{h\in \mathbf{Z}\\ |h|\leqslant \mathrm{N}}}\widehat{f} (h)e^{2i\pi ht}\Bigl (1 - \frac{|h|}{\mathrm{N}}\Bigr)
\]

for \(t \in T\) . Let \(X \subset T\) a set such that the restriction of f to X is continuous. Prove that \(f_{N}\) converges to f uniformly on X.

\begin{enumerate}
\def\labelenumi{\alph{enumi})}
\setcounter{enumi}{2}
\tightlist
\item
  Let \(X \subset T\) an open set such that the restriction of f to X is of class \(C^{1}\) . Prove that the Fourier series of f converges symmetrically to f uniformly on X.
\end{enumerate}

\begin{enumerate}
\def\labelenumi{\arabic{enumi})}
\setcounter{enumi}{53}
\tightlist
\item
  Suppose that G is not compact. Let \(\widehat{G}_{d}\) the dual group \(\widehat{G}\) endowed with the discrete topology and \(\widetilde{G}\) the dual group of \(\widehat{G}_{d}\) ; this is a compact topological group.
\end{enumerate}

\begin{enumerate}
\def\labelenumi{\alph{enumi})}
\tightlist
\item
  The map which associates to \(x \in \mathrm{G}\) the character \(\chi \mapsto \chi(x)\) is a continuous injective homomorphism of \(\mathrm{G}\) into \(\widetilde{\mathrm{G}}\) ; its image is dense in \(\widetilde{\mathrm{G}}\) .
\end{enumerate}

From now on we identify G with a subset of \(\widetilde{G}\) .

\begin{enumerate}
\def\labelenumi{\alph{enumi})}
\setcounter{enumi}{1}
\item
  Let K be a compact topological group and \(\varphi\colon\mathrm{G}\to\mathrm{K}\) a continuous homomorphism. There exists a unique homomorphism \(\widetilde{\varphi}\) of \(\widetilde{\mathrm{G}}\) into K which extends \(\varphi\) .
\item
  Let \(f \in \mathcal{C}_b(\mathrm{G})\) . The following conditions are equivalent:
\end{enumerate}

\begin{enumerate}
\def\labelenumi{(\roman{enumi})}
\item
  There exists a function \(\widetilde{f} \in \mathcal{C}(\widetilde{\mathrm{G}})\) which extends \(f\) ;
\item
  The function \(f\) belongs to the closure in \(\mathcal{C}(\mathrm{G})\) of the subspace of \(\mathcal{C}(\mathrm{G})\) generated by the characters \(\chi\in\widehat{\mathrm{G}}\) ;
\item
  The set of \(\gamma(x)f\) for \(x\in G\) is precompact in \(\mathcal{C}(G)\) .
\end{enumerate}

(To verify that (iii) implies (i), show that \(f\) is uniformly continuous. Then consider the closure of the \(\pmb{\gamma}(x)f\) ; this is a compact metric space K. Prove that \(\varrho \colon x \mapsto \pmb{\gamma}(x)\) is a continuous homomorphism of G into the group of isometries of K; then apply question \(b)\) to extend \(\varrho\) to \(\widetilde{\varrho}\) and set \(\widetilde{f}(x) = (\widetilde{\varrho}(x)f)(e).\)

One says that a function \(f \in \mathcal{C}_b(\mathrm{G})\) satisfying these equivalent conditions is almost periodic.

\begin{enumerate}
\def\labelenumi{\alph{enumi})}
\setcounter{enumi}{3}
\item
  Let f be an almost periodic function on G. For every \(\varepsilon > 0\) and for every compact subset F of G, there exists \(x \in G - F\) such that \(\|\boldsymbol{\gamma}(x)f - f\| \leqslant \varepsilon\) .
\item
  The set of almost periodic functions on G is a star subalgebra of \(\mathcal{C}_{b}(\mathrm{G})\) .
\end{enumerate}

\begin{enumerate}
\def\labelenumi{\arabic{enumi})}
\setcounter{enumi}{54}
\tightlist
\item
  We keep the notation from the previous exercise. We endow \(\widetilde{G}\) of the normalized Haar measure \(\nu\) Let f be an almost periodic function on G and \(\widetilde{f}\)
\end{enumerate}

its extension to \(\widetilde{\mathrm{G}}\) . We set

\[
\mathrm{A} (f) = \int_ {\widetilde {\mathrm{G}}} \widetilde {f} (x) \nu (x).
\]

\begin{enumerate}
\def\labelenumi{\alph{enumi})}
\tightlist
\item
  Let \(\chi \in \widehat{\mathbf{G}}\) . The function \(\overline{\chi} f\) is almost periodic. We then define \(\mathcal{P}(f)(\chi) = \mathrm{A}(\overline{\chi} f)\) . We have
\end{enumerate}

\[
\mathrm{A} (| f | ^ {2}) = \sum_ {\chi \in \widehat {\mathrm{G}}} | \mathcal {P} (f) (\chi) | ^ {2}
\]

    and in particular \(A(\overline{\chi}f)=0\) except for a countable set of \(\chi\) b) Let f be an almost periodic function on R. For every \(y \in R\) , one has

\[
\mathcal {P} (f) (y) = \lim _ {\mathrm{T} \to + \infty} \frac {1}{2 \mathrm{T}} \int_ {- \mathrm{T}} ^ {\mathrm{T}} f (x) e ^ {- 2 i \pi x y} d x.
\]

\begin{enumerate}
\def\labelenumi{\alph{enumi})}
\setcounter{enumi}{2}
\tightlist
\item
  Let \(\Lambda\) a finite subset of \(\mathbf{R}\) . The function defined on \(\mathbf{R}\) by
\end{enumerate}

\[
f (x) = \sum_ {\lambda \in \Lambda} e ^ {i \lambda x}
\]

is almost periodic on \(\mathbf{R}\) to determine as a function of \(\Lambda\) the set of \(y \in \mathbf{R}\) such that \(\mathcal{P}(f)(y) \neq 0\) .

\begin{enumerate}
\def\labelenumi{\arabic{enumi})}
\setcounter{enumi}{55}
\tightlist
\item
  We set
\end{enumerate}

\[
k (z) = \left(\frac {\sin (\pi z)}{\pi z}\right) ^ {2}
\]

if \(z \in \mathbf{C} - \mathbf{Z}\) and \(k(z) = 1\) if \(z \in \mathbf{Z}\) . The function \(k\) belongs to the space \(\mathcal{E}_{1/2}\) (Exercise 36).

\begin{enumerate}
\def\labelenumi{\alph{enumi})}
\tightlist
\item
  For every complex number \(z \in \mathbf{C} - \mathbf{Z}\) , one has
\end{enumerate}

\[
\sum_ {n \in {\bf Z}} \frac {1}{(z - n) ^ {2}} = \left(\frac {\pi}{\sin (\pi z)}\right) ^ {2}
\]

(apply the Poisson formula).

\begin{enumerate}
\def\labelenumi{\alph{enumi})}
\setcounter{enumi}{1}
\tightlist
\item
  For \(z \in \mathbf{C} - \mathbf{Z}\) , we define
\end{enumerate}

\[
h (z) = \Big (\frac {\sin (\pi z)}{\pi} \Big) ^ {2} \Big \{\sum_ {n \in \mathbf {Z}} \frac {\mathrm{s} (n)}{(z - n) ^ {2}} + \frac {2}{z} \Big \},
\]

where \(s(n)\) denotes the sign function. The function h extends to an entire function belonging to the space \(E_{2}\) . We set \(b = h + k\) .

\begin{enumerate}
\def\labelenumi{\alph{enumi})}
\setcounter{enumi}{2}
\item
  For every \(x \in \mathbf{R}\) , one has \(1 - k(x) \leqslant h(x) \leqslant 1\) .
\item
  For every \(x \in R\) , one has \(|\mathrm{s}(x) - h(x)| \leqslant k(x)\) .
\item
    We have \(b \geqslant s\) ; the function b-s is integrable on R and satisfies \(\int(b - s) = 1\) . f) Let \(\alpha < \beta\) be real numbers and \(\varphi\) the characteristic function of the interval \([\alpha, \beta]\) . The function \(f_{\alpha,\beta}(z) = \frac{1}{2}(b(\beta - z) + b(z - \alpha))\) satisfies \(\varphi \leqslant f_{\alpha,\beta}\) and \(\int_{\mathbf{R}}(f_{\alpha,\beta} - \varphi) = 1\) .
\end{enumerate}

\(g)\) *There exists an entire function \(\sigma\) such that \(f_{\alpha,\beta}(x)=|\sigma(x)|^{2}\) for all \(x\in\mathbf{R}\) . The Fourier transform of the restriction of \(\sigma\) to \(\mathbf{R}\) is supported in \([-1/2,1/2].*\)

\begin{enumerate}
\def\labelenumi{\arabic{enumi})}
\setcounter{enumi}{56}
\tightlist
\item
  We keep the notations of the previous exercise. Let M, \(N \geqslant 1\) be integers and let \((a_{n})_{\mathrm{M} < n \leqslant \mathrm{M} + \mathrm{N}}\) be a family of complex numbers. Let \(\delta > 0\) be a real number, \(R \geqslant 1\) be an integer and \((\xi_{r})_{1 \leqslant r \leqslant \mathrm{R}}\) a family of real numbers such that \(\inf_{h \in \mathbf{Z}} |\xi_{r} - \xi_{r'} - h| \geqslant \delta\) for all \(r \neq r'\) .
\end{enumerate}

For every \(x\in \mathbf{R}\) , set

\[
\mathrm{S} (x) = \sum_ {n = \mathrm{M} + 1} ^ {\mathrm{M} + \mathrm{N}} a _ {n} e ^ {2 i \pi n x}.
\]

\begin{enumerate}
\def\labelenumi{\alph{enumi})}
\tightlist
\item
  We have
\end{enumerate}

\[
\sum_ {r = 1} ^ {\mathrm{R}} | \mathrm{S} (\xi_ {r}) | ^ {2} \leqslant (\mathrm{N} + \delta^ {- 1} - 1) \sum_ {n = \mathrm{M} + 1} ^ {\mathrm{M} + \mathrm{N}} | a _ {n} | ^ {2}
\]

("large sieve inequality").

Consider the function \(\tau\colon z\mapsto\sigma(\delta z)\) , where \(\sigma\) is the function of the questions \(f\) ) and \(g\) ) of the preceding exercise for the interval \([\alpha,\beta]=[ \delta(M+1),\delta(M+N)]\) ; set

\[
\mathrm{S} ^ {*} (x) = \sum_ {n = \mathrm{M} + 1} ^ {\mathrm{M} + \mathrm{N}} a (n) \tau (n) ^ {- 1} e ^ {2 i \pi n x}
\]

and observe that

\[
\mathrm{S} (x) = \int_ {- \delta / 2} ^ {\delta / 2} \widehat {\tau} (t) \mathrm{S} ^ {*} (t + x) d t.
\]

Set \(x = \xi_r\) , apply the Cauchy-Schwarz inequality and sum over \(r\) .)

\begin{enumerate}
\def\labelenumi{\alph{enumi})}
\setcounter{enumi}{1}
\tightlist
\item
  The constant \(N + \delta^{-1} - 1\) is optimal.
\end{enumerate}

\begin{enumerate}
\def\labelenumi{\arabic{enumi})}
\setcounter{enumi}{57}
\tightlist
\item
  For every nonzero real number \(\lambda\) , we denote by \(\mu_{\lambda}\) the discrete probability measure on \(\mathbf{R}\) such that \(\mu_{\lambda}(\lambda) = \mu_{\lambda}(-\lambda) = \frac{1}{2}\) .
\end{enumerate}

We fix \(\lambda \in ]0,1[\) .

\begin{enumerate}
\def\labelenumi{\alph{enumi})}
\tightlist
\item
  Let us set \(\nu_0 = \mu_1\) and \(\nu_{n+1} = \mu_{\lambda^{n+1}} * \nu_n\) for every integer \(n \geqslant 0\) . The sequence \((\nu_n)_{n \in \mathbf{N}}\) converges to a probability measure \(\nu_\lambda\) on \(\mathbf{R}\) The characteristic function of \(\nu_\lambda\) is the function on \(\mathbf{R}\) defined by
\end{enumerate}

\[
\varphi_ {\lambda} (x) = \prod_ {n = 0} ^ {+ \infty} \cos (\lambda^ {n} x).
\]

\begin{enumerate}
\def\labelenumi{\alph{enumi})}
\setcounter{enumi}{1}
\tightlist
\item
  The measure \(\nu_{\lambda}\) is the unique positive measure \(\nu\) of total mass 1 on R such that
\end{enumerate}

\[
\nu = \frac {1}{2} \ell_ {1} (\nu) + \frac {1}{2} \ell_ {2} (\nu),
\]

where \(\ell_{1}\) (resp. \(\ell_{2}\) ) is the map \(x \mapsto \lambda x - 1\) (resp. \(x \mapsto \lambda x + 1\) ) from R to R.

\begin{enumerate}
\def\labelenumi{\alph{enumi})}
\setcounter{enumi}{2}
\item
  If the measure \(\nu_{\lambda}\) is not a base measure, the Lebesgue measure, then \(\nu_{\lambda}\) is singular with respect to the Lebesgue measure. (Consider the Lebesgue decomposition of \(\nu_{\lambda}\) with respect to the Lebesgue measure, cf.~INT, V, §5, no. 7, th. 3.)
\item
  If \(0 < \lambda < \frac{1}{2}\) , then the measure \(\nu_{\lambda}\) is singular with respect to the Lebesgue measure.
\item
  If \(\lambda = \frac{1}{2}\) , then the measure \(\nu_{\lambda}\) is equal to the uniform measure on {[}-2, 2{]}. If \(\lambda > 1/2\) , the support of \(\nu_{\lambda}\) is the interval \([-(1-\lambda)^{-1}, (1-\lambda)^{-1}]\) .
\item
  A Pisot number is a real algebraic integer \(x\) such that all the conjugates \(y \in \mathbf{C}\) of \(x\) , other than \(x\) , satisfy \(|y| < 1\) . If \(\lambda^{-1}\) is a Pisot number, then the function \(\varphi_{\lambda}\) does not tend to 0 at infinity. In particular, the measure \(\nu_{\lambda}\) is singular with respect to the Lebesgue measure.
\item
  There exists \(\lambda > 1/2\) such that \(\nu_{\lambda}\) is singular with respect to the Lebesgue measure.
\end{enumerate}
